\documentclass[12pt]{article}

% === CÀI ĐẶT TRANG ===
\usepackage[margin=2cm]{geometry}
\usepackage[utf8]{inputenc}
\usepackage[T5]{fontenc}
\usepackage[vietnamese]{babel}

% === TOÁN HỌC & HÌNH VẼ ===
\usepackage{amsmath, amssymb, amsthm}
\usepackage{graphicx}
\usepackage{float}
\usepackage{tikz}
\usepackage{pgfplots}
\pgfplotsset{compat=1.18}

% === ĐỊNH DẠNG ===
\usepackage{enumitem}
\usepackage{multicol}
\usepackage{fancyhdr}
\usepackage[hidelinks]{hyperref}

% Sửa lỗi headheight quá nhỏ
\setlength{\headheight}{14pt}

% Thiết lập khoảng cách cho list để văn bản thoáng hơn
\setlist[enumerate]{itemsep=4pt, topsep=4pt}

% === HEADER/FOOTER ===
\pagestyle{fancy}
\fancyhf{}
\fancyhead[L]{\small Đề cương ôn tập KTTX 2}
\fancyhead[R]{\small Lớp 2026 - L01}
\fancyfoot[C]{Trang \thepage}
\renewcommand{\headrulewidth}{0.4pt}

\begin{document}
	
	% ================================================
	% TRANG BÌA
	% ================================================
	\begin{titlepage}
		\centering
		\vspace*{4cm}
		
		{\Huge \textbf{ĐỀ CƯƠNG ÔN TẬP}}\\[1cm]
		{\LARGE \textbf{KIỂM TRA THƯỜNG XUYÊN 2}}\\[1.5cm]
		
		{\Large \textbf{Lớp: 2026 - L01}}\\[0.5cm]
		
		{\large \textit{(Lưu hành nội bộ)}}
		
		\vfill
	\end{titlepage}
	
	% ================================================
	% ĐỀ SỐ 1
	% ================================================
	\newpage
	\begin{center}
		{\LARGE \textbf{ĐỀ 1}}\\[4pt]
	\end{center}
	\vspace{4mm}
	
	% --- PHẦN 1: TRẮC NGHIỆM ---
	\noindent\textbf{PHẦN I. Câu trắc nghiệm nhiều phương án lựa chọn.} Thí sinh trả lời từ câu 1 đến câu 12. Mỗi câu hỏi thí sinh chỉ chọn một phương án.
	\vspace{2mm}
	
	\begin{enumerate}[label=\textbf{Câu \arabic*.}, leftmargin=*]
		
		% CÂU 1: Vector trong không gian (Hình hộp chữ nhật)
		\item Cho hình hộp chữ nhật $ABCD.A'B'C'D'$. Khẳng định nào sau đây là đúng?
		\begin{center}
			\begin{tikzpicture}[scale=0.8, line join=round, line cap=round]
				% Định nghĩa các đỉnh
				\coordinate (A) at (0,0);
				\coordinate (B) at (4,0);
				\coordinate (C) at (5.5,1.5);
				\coordinate (D) at (1.5,1.5);
				\coordinate (A') at (0,4);
				\coordinate (B') at (4,4);
				\coordinate (C') at (5.5,5.5);
				\coordinate (D') at (1.5,5.5);
				
				% Vẽ các nét khuất (đứt)
				\draw[dashed, thick] (A) -- (D) (D) -- (C) (D) -- (D');
				% Vẽ các nét thấy (liền)
				\draw[thick] (A) -- (B) -- (C) -- (C') -- (B') -- (A') -- cycle;
				\draw[thick] (A) -- (A') (B) -- (B') (B') -- (C');
				\draw[thick] (D') -- (A') (D') -- (C');
				
				% Gắn nhãn
				\node[below left] at (A) {$A$};
				\node[below right] at (B) {$B$};
				\node[right] at (C) {$C$};
				\node[above left] at (D) {$D$};
				\node[above left] at (A') {$A'$};
				\node[above left] at (B') {$B'$};
				\node[right] at (C') {$C'$};
				\node[above left] at (D') {$D'$};
			\end{tikzpicture}
		\end{center}
		\begin{center}
			\begin{tabular}{p{0.42\textwidth} p{0.42\textwidth}}
				A. $\overrightarrow{CB} + \overrightarrow{CD} + \overrightarrow{CC'} = \overrightarrow{CA'}$. &
				B. $\overrightarrow{CB} + \overrightarrow{CD} + \overrightarrow{CC'} = \overrightarrow{AC'}$. \\[2mm]
				C. $\overrightarrow{CB} + \overrightarrow{CD} + \overrightarrow{CC'} = \overrightarrow{BD'}$. &
				D. $\overrightarrow{CB} + \overrightarrow{CD} + \overrightarrow{CC'} = \overrightarrow{DB'}$.
			\end{tabular}
		\end{center}
		
		% CÂU 2: Max/Min hàm phân thức
		\item Gọi $M, m$ lần lượt là giá trị lớn nhất và giá trị nhỏ nhất của hàm số $y = \dfrac{x+2}{x-1}$ trên đoạn $[2; 4]$. Tính giá trị của biểu thức $T = M + m$.
		\begin{multicols}{4}
			\begin{enumerate}[label=\Alph*.]
				\item $T = 6$.
				\item $T = 2$.
				\item $T = 4$.
				\item $T = \dfrac{10}{3}$.
			\end{enumerate}
		\end{multicols}
		
		% CÂU 3: Đọc đồ thị tìm khoảng đơn điệu
		\item Cho hàm số $y = f(x)$ có đồ thị như hình vẽ bên dưới. Hàm số đã cho nghịch biến trên khoảng nào dưới đây?
		\begin{center}
			\begin{tikzpicture}
				\begin{axis}[
					axis lines = middle,
					xlabel = {$x$}, ylabel = {$y$},
					ymin = -2, ymax = 4.5,
					xmin = -2.5, xmax = 2.5,
					xtick = {-1, 1}, ytick = {-1, 1, 3},
					ticklabel style = {fill=white, inner sep=1pt},
					width=8cm, height=6.5cm,
					samples=100,
					axis line style={->}
					]
					% Đồ thị y = x^3 - 3x + 1
					\addplot [domain=-2.1:2.1, thick, blue] {x^3 - 3*x + 1};
					
					% Gióng cực đại
					\draw[dashed, gray] (-1,0) -- (-1,3) -- (0,3);
					\fill (-1,3) circle (1.5pt);
					% Gióng cực tiểu
					\draw[dashed, gray] (1,0) -- (1,-1) -- (0,-1);
					\fill (1,-1) circle (1.5pt);
					
					% Gốc tọa độ
					\node[below left] at (axis cs:0,0) {$O$};
				\end{axis}
			\end{tikzpicture}
		\end{center}
		\begin{multicols}{4}
			\begin{enumerate}[label=\Alph*.]
				\item $(-1; 1)$.
				\item $(-\infty; -1)$.
				\item $(1; +\infty)$.
				\item $(-1; 3)$.
			\end{enumerate}
		\end{multicols}
		
		% CÂU 4: Thống kê (Khoảng tứ phân vị)
		\item Số lượng sản phẩm hoàn thành trong một ngày của một nhóm công nhân được thống kê bởi bảng ghép nhóm sau:
		\begin{center}
			\begin{tabular}{|c|c|}
				\hline
				\textbf{Số lượng sản phẩm} & \textbf{Số công nhân} \\
				\hline
				$[10; 20)$ & 5 \\
				\hline
				$[20; 30)$ & 12 \\
				\hline
				$[30; 40)$ & 18 \\
				\hline
				$[40; 50)$ & 10 \\
				\hline
				$[50; 60)$ & 5 \\
				\hline
			\end{tabular}
		\end{center}
		Hãy tìm khoảng tứ phân vị $\Delta_Q$ của mẫu số liệu ghép nhóm cho bởi bảng trên.
		\begin{multicols}{4}
			\begin{enumerate}[label=\Alph*.]
				\item $\Delta_Q = 16,25$.
				\item $\Delta_Q = 15,75$.
				\item $\Delta_Q = 14,5$.
				\item $\Delta_Q = 16,5$.
			\end{enumerate}
		\end{multicols}
		
		% CÂU 5: Tìm số tiệm cận đứng
		\item Tìm số tiệm cận đứng của đồ thị hàm số $y = \dfrac{x^2 - x - 6}{x^2 - 9}$.
		\begin{multicols}{4}
			\begin{enumerate}[label=\Alph*.]
				\item $1$.
				\item $2$.
				\item $0$.
				\item $3$.
			\end{enumerate}
		\end{multicols}
		
		% CÂU 6: Nhận diện dấu hệ số hàm bậc 3 từ đồ thị
		\item Cho hàm số $y = ax^3 + bx^2 + cx + d$ $(a \neq 0)$ có đồ thị như hình vẽ bên. Mệnh đề nào dưới đây đúng?
		\begin{center}
			\begin{tikzpicture}
				\begin{axis}[
					axis lines = middle,
					xlabel = {$x$}, ylabel = {$y$},
					ymin = -3.5, ymax = 4.5,
					xmin = -2.5, xmax = 2.5,
					xtick = {-1, 1, 2}, ytick = {1},
					ticklabel style = {fill=white, inner sep=1pt},
					width=8cm, height=6.5cm,
					samples=100,
					axis line style={->}
					]
					% Đồ thị y = x^3 + 1.5x^2 - 2x + 1 
					% Hàm này đi lên-xuống-lên, cắt tung độ tại dương, cực đại hoành âm (trị tuyệt đối lớn hơn), cực tiểu hoành dương.
					\addplot [domain=-2.5:1.5, thick, blue] {x^3 + 1.5*x^2 - 2*x + 1};
					
					% Gốc tọa độ
					\node[below left] at (axis cs:0,0) {$O$};
				\end{axis}
			\end{tikzpicture}
		\end{center}
		\begin{center}
			\begin{tabular}{p{0.42\textwidth} p{0.42\textwidth}}
				A. $a > 0; b > 0; c < 0; d > 0$. &
				B. $a > 0; b < 0; c < 0; d > 0$. \\[2mm]
				C. $a < 0; b > 0; c > 0; d < 0$. &
				D. $a > 0; b > 0; c > 0; d > 0$.
			\end{tabular}
		\end{center}
		
		% ======================================================
		% TIẾP TỤC ĐỀ 1 - TỪ CÂU 7 ĐẾN CÂU 12
		% ======================================================
		
		% CÂU 7: Tích vô hướng của hai vecto (Dựa trên câu 10 ảnh cuối)
		\item Trong không gian $Oxyz$, cho $\vec{a} = 2\vec{i} - \vec{j} + 3\vec{k}$ và $\vec{b} = (-1; 4; 2)$. Tính $\vec{a} \cdot \vec{b}$.
		\begin{multicols}{4}
			\begin{enumerate}[label=\Alph*.]
				\item $\vec{a} \cdot \vec{b} = 0$.
				\item $\vec{a} \cdot \vec{b} = 1$.
				\item $\vec{a} \cdot \vec{b} = (1; 3; 5)$.
				\item $\vec{a} \cdot \vec{b} = -2$.
			\end{enumerate}
		\end{multicols}
		
		% CÂU 8: Đọc bảng biến thiên tìm Max/Min (Dựa trên câu 2 ảnh giữa)
		\item Cho hàm số $y = f(x)$ có bảng biến thiên như sau:
		\begin{center}
			\begin{tikzpicture}[xscale=1.2, yscale=0.8]
				% Khung bảng
				\draw (0,0) rectangle (8,3);
				\draw (0,1) -- (8,1);
				\draw (0,2) -- (8,2);
				\draw (1,0) -- (1,3);
				
				% Dòng x
				\node at (0.5, 2.5) {$x$};
				\node at (1.5, 2.5) {$-\infty$};
				\node at (3.5, 2.5) {$-1$};
				\node at (5.5, 2.5) {$3$};
				\node at (7.5, 2.5) {$+\infty$};
				
				% Dòng f'(x)
				\node at (0.5, 1.5) {$f'(x)$};
				\node at (2.5, 1.5) {$+$};
				\node at (3.5, 1.5) {$0$};
				\node at (4.5, 1.5) {$-$};
				\node at (5.5, 1.5) {$0$};
				\node at (6.5, 1.5) {$+$};
				
				% Dòng f(x)
				\node at (0.5, 0.5) {$f(x)$};
				\node at (1.5, 0.2) {$-\infty$};
				\node at (3.5, 0.7) {$4$};
				\node at (5.5, 0.2) {$-2$};
				\node at (7.5, 0.7) {$+\infty$};
				
				% Mũi tên
				\draw[->, thick, blue] (1.8, 0.3) -- (3.2, 0.6);
				\draw[->, thick, blue] (3.8, 0.6) -- (5.2, 0.3);
				\draw[->, thick, blue] (5.8, 0.3) -- (7.2, 0.6);
			\end{tikzpicture}
		\end{center}
		Trên khoảng $(-\infty; 2)$, giá trị lớn nhất của hàm số đã cho bằng:
		\begin{multicols}{4}
			\begin{enumerate}[label=\Alph*.]
				\item $3$.
				\item $-2$.
				\item $4$.
				\item $-1$.
			\end{enumerate}
		\end{multicols}
		
		% CÂU 9: Tọa độ điểm thỏa mãn biểu thức vecto (Dựa trên câu 10 ảnh giữa)
		\item Trong không gian với hệ trục tọa độ $Oxyz$, cho vecto $\vec{v} = (-2; 3; 1)$ và điểm $P(3; -1; 2)$. Tọa độ điểm $Q$ thỏa mãn $\overrightarrow{PQ} = \vec{v}$ là:
		\begin{multicols}{4}
			\begin{enumerate}[label=\Alph*.]
				\item $Q(-5; 4; -1)$.
				\item $Q(1; 2; 3)$.
				\item $Q(5; -4; 1)$.
				\item $Q(-1; -2; -3)$.
			\end{enumerate}
		\end{multicols}
		
		% CÂU 10: Thống kê tính Độ lệch chuẩn (Dựa trên câu 12 ảnh đầu)
		\item Để chuẩn bị cho giải điền kinh cấp trường, bạn Hùng đã tự ghi lại cự ly chạy (đơn vị: km) của mình trong 20 ngày luyện tập liên tiếp và thu được kết quả ở bảng sau:
		\begin{center}
			\begin{tabular}{|c|c|c|c|c|c|}
				\hline
				\textbf{Cự ly chạy (km)} & $[0; 2)$ & $[2; 4)$ & $[4; 6)$ & $[6; 8)$ & $[8; 10)$ \\
				\hline
				\textbf{Số ngày} & 2 & 4 & 8 & 4 & 2 \\
				\hline
			\end{tabular}
		\end{center}
		Độ lệch chuẩn của mẫu số liệu ghép nhóm trên có giá trị gần nhất với giá trị nào dưới đây?
		\begin{multicols}{4}
			\begin{enumerate}[label=\Alph*.]
				\item $4,80$.
				\item $2,25$.
				\item $2,19$.
				\item $5,00$.
			\end{enumerate}
		\end{multicols}
		
		% CÂU 11: Tính tọa độ vecto cơ bản (Dựa trên câu 6 ảnh đầu)
		\item Trong không gian $Oxyz$, cho $\overrightarrow{OA} = -\vec{i} + 3\vec{j} + 4\vec{k}$ và $\overrightarrow{OB} = 5\vec{i} - \vec{k}$. Tìm tọa độ của vecto $\overrightarrow{AB}$.
		\begin{center}
			\begin{tabular}{p{0.42\textwidth} p{0.42\textwidth}}
				A. $\overrightarrow{AB} = (6; -3; -5)$. &
				B. $\overrightarrow{AB} = (-6; 3; 5)$. \\[2mm]
				C. $\overrightarrow{AB} = (4; 3; 3)$. &
				D. $\overrightarrow{AB} = (6; 3; -5)$.
			\end{tabular}
		\end{center}
		
		% CÂU 12: Phép cộng vecto hình không gian (Dựa trên câu 9 ảnh giữa)
		\item Cho hình lập phương $ABCD.A'B'C'D'$ (như hình minh họa bên dưới) có độ dài mỗi cạnh bằng $1$. Tính độ dài của vecto $\overrightarrow{AB} + \overrightarrow{B'C'}$.
		\begin{center}
			\begin{tikzpicture}[scale=0.9, line join=round, line cap=round]
				% Tọa độ các đỉnh của hình lập phương
				\coordinate (D) at (0,0);
				\coordinate (C) at (3,0);
				\coordinate (B) at (4,1.2);
				\coordinate (A) at (1,1.2);
				\coordinate (D') at (0,3);
				\coordinate (C') at (3,3);
				\coordinate (B') at (4,4.2);
				\coordinate (A') at (1,4.2);
				
				% Vẽ nét khuất
				\draw[dashed, thick] (A) -- (D) (A) -- (B) (A) -- (A') (A) -- (C);
				% Vẽ nét liền
				\draw[thick] (D) -- (C) -- (B) -- (B') -- (C') -- (D') -- cycle;
				\draw[thick] (D) -- (D') (C) -- (C') (A') -- (B') (A') -- (C') (A') -- (D');
				
				% Nhãn
				\node[below left] at (D) {$D$};
				\node[below right] at (C) {$C$};
				\node[right] at (B) {$B$};
				\node[above left] at (A) {$A$};
				\node[left] at (D') {$D'$};
				\node[right] at (C') {$C'$};
				\node[above right] at (B') {$B'$};
				\node[above left] at (A') {$A'$};
			\end{tikzpicture}
		\end{center}
		\begin{multicols}{4}
			\begin{enumerate}[label=\Alph*.]
				\item $\sqrt{3}$.
				\item $1$.
				\item $\sqrt{2}$.
				\item $2$.
			\end{enumerate}
		\end{multicols}
		
	\end{enumerate}
	
	\vspace{4mm}
	
	
	% ======================================================
	% PHẦN 2: TRẮC NGHIỆM ĐÚNG SAI
	% ======================================================
	\noindent\textbf{PHẦN II. Câu trắc nghiệm đúng sai.} Thí sinh trả lời từ câu 1 đến câu 2. Trong mỗi ý a), b), c), d) ở mỗi câu, thí sinh chọn đúng hoặc sai.
	\vspace{2mm}
	
	\begin{enumerate}[label=\textbf{Câu \arabic*.}, leftmargin=*]
	
	% CÂU 1 ĐÚNG SAI: Khảo sát tính chất hàm phân thức bậc 2/bậc 1
	\item Cho hàm số $y = f(x) = \dfrac{x^2 - x + 2}{x - 2}$ có đồ thị là đường cong $(C)$. Xét tính đúng hay sai của các mệnh đề sau:
	\begin{enumerate}[label=\alph*)]
		\item Hàm số $f(x)$ có tập xác định là $\mathbb{R} \setminus \{2\}$.
		\item Đồ thị $(C)$ có đường tiệm cận đứng là đường thẳng $x = -2$.
		\item Tâm đối xứng của đồ thị hàm số là điểm $I(2; 3)$.
		\item Trên đồ thị $(C)$ tồn tại đúng $4$ điểm có tọa độ là các số nguyên.
	\end{enumerate}
	
	\vspace{4mm}
	% CÂU 2 ĐÚNG SAI: Tọa độ điểm và vecto trong không gian Oxyz
	\item Trong không gian $Oxyz$, cho các điểm $A(3; -1; 2)$, $B(1; 1; -2)$ và $C(0; 2; -3)$. Xét tính đúng hay sai của các mệnh đề sau:
	\begin{enumerate}[label=\alph*)]
		\item Tọa độ của vecto $\overrightarrow{AB}$ là $(-2; 2; -4)$.
		\item Độ dài của vecto $\overrightarrow{AB}$ bằng $2\sqrt{5}$.
		\item Điểm $M$ thỏa mãn hệ thức $\overrightarrow{AB} + \overrightarrow{CM} = \vec{0}$ có tọa độ là $M(2; 0; 1)$.
		\item Gọi $N$ là giao điểm của đường thẳng đi qua hai điểm $A, B$ với mặt phẳng tọa độ $(Oxy)$. Tọa độ của điểm $N$ là $N(2; 0; 0)$.
	\end{enumerate}
	
	\end{enumerate} % Kết thúc Phần II
	
	% ======================================================
	% PHẦN 3: TRẮC NGHIỆM TRẢ LỜI NGẮN
	% ======================================================
	\vspace{4mm}
	\noindent\textbf{PHẦN III. Câu trắc nghiệm trả lời ngắn.}
	\vspace{2mm}
	
	\begin{enumerate}[label=\textbf{Câu \arabic*.}, leftmargin=*]
		
		% CÂU 1 TRẢ LỜI NGẮN: Tổng cực đại, cực tiểu hàm phân thức
		\item Cho hàm số $y = f(x) = \dfrac{x^2 - 2x + 5}{x - 1}$. Tổng giá trị cực đại và giá trị cực tiểu của hàm số đã cho bằng bao nhiêu?
		
		\vspace{2mm}
		% CÂU 2 TRẢ LỜI NGẮN: Thống kê (Hiệu số 2 khoảng tứ phân vị)
		\item Điều tra về thời gian tự học ở nhà (giờ/tuần) của học sinh ở hai lớp 11A và 11B tại một trường THPT, ta có hai mẫu số liệu ghép nhóm sau:
		
		\begin{center}
			\renewcommand{\arraystretch}{1.2}
			\begin{tabular}{|c|c|c|c|c|c|}
				\hline
				\textbf{Thời gian (giờ)} & $[5; 10)$ & $[10; 15)$ & $[15; 20)$ & $[20; 25)$ & $[25; 30)$ \\
				\hline
				\textbf{Số học sinh 11A} & 2 & 10 & 18 & 8 & 2 \\
				\hline
				\textbf{Số học sinh 11B} & 4 & 12 & 14 & 8 & 2 \\
				\hline
			\end{tabular}
		\end{center}
		Tính hiệu số giữa khoảng tứ phân vị của mẫu số liệu ghép nhóm lớp 11B và khoảng tứ phân vị của mẫu số liệu ghép nhóm lớp 11A.
		
		\vspace{2mm}
		% CÂU 3 TRẢ LỜI NGẮN: Tính góc giữa 2 vecto dựa vào độ dài
		\item Cho 2 vecto $\vec{u}$ và $\vec{v}$ thỏa mãn $|\vec{u}| = 2$, $|\vec{v}| = 3$ và $|\vec{u} - \sqrt{3}\vec{v}| = \sqrt{13}$. Tính góc $(\vec{u}, \vec{v})$ (viết kết quả tính theo đơn vị độ).
		
		\vspace{2mm}
		% CÂU 4 TRẢ LỜI NGẮN: Hình hộp trong Oxyz
		\item Trong không gian $Oxyz$, cho hình hộp $ABCD.A'B'C'D'$ có tọa độ các điểm $B(3; 1; 0)$, $D(-1; 3; 2)$, $A'(4; 3; 5)$ và $C'(2; -1; 3)$. Giả sử tọa độ điểm $B'(x; y; z)$, tính giá trị của biểu thức $T = x + y + z$.
		
	\end{enumerate}
	
	
	% ======================================================
	% PHẦN 4: TỰ LUẬN
	% ======================================================
	\vspace{4mm}
	\noindent\textbf{PHẦN IV. Tự luận.} Thí sinh trình bày bài giải chi tiết từ bài 1 đến bài 2.
	\vspace{2mm}
	
	\begin{enumerate}[label=\textbf{Bài \arabic*.}, leftmargin=*]
		
		% BÀI 1: KHẢO SÁT HÀM SỐ (1.5 ĐIỂM)
		\item  Cho hàm số $y = \dfrac{x^2 - 3x + 3}{x - 1}$ có đồ thị là đường cong $(C)$.
		\begin{enumerate}[label=\alph*)]
			\item Khảo sát sự biến thiên (đồng biến, nghịch biến, cực trị, tiệm cận) và xác định tọa độ tâm đối xứng của đồ thị hàm số đã cho.
			\item Viết phương trình tiếp tuyến của đồ thị $(C)$, biết tiếp tuyến đó song song với đường thẳng $d: y = -3x + 2025$.
		\end{enumerate}
		
		\vspace{4mm}
		% BÀI 2: VECTO TRONG KHÔNG GIAN (1.5 ĐIỂM)
		\item  
		\begin{enumerate}[label=\alph*)]
			\item Cho hình chóp $S.ABCD$ có đáy $ABCD$ là hình bình hành. Gọi $M, N$ lần lượt là trung điểm của các cạnh $BC$ và $AD$. Chứng minh đẳng thức vecto sau:
			\[ \overrightarrow{SM} + \overrightarrow{SN} = \overrightarrow{SB} + \overrightarrow{SD} \]
			\item Cho hình chóp tứ giác đều $S.ABCD$ có tất cả các cạnh đều bằng $a$ (Hình minh họa bên cạnh). Tính các tích vô hướng sau theo $a$: 
			\[ \overrightarrow{SA} \cdot \overrightarrow{SC} \quad \text{và} \quad \overrightarrow{SA} \cdot \overrightarrow{BC} \]
		\end{enumerate}
		\hfill
		\begin{minipage}[t]{0.3\textwidth}
			\vspace{-1.5cm} % Đẩy hình lên cho ngang hàng với text
			\begin{center}
				\begin{tikzpicture}[scale=0.9, line join=round, line cap=round]
					% Định nghĩa các đỉnh của chóp tứ giác đều
					\coordinate (A) at (0,0);
					\coordinate (D) at (3.5,0);
					\coordinate (B) at (-1.5,-1.5);
					\coordinate (C) at (2,-1.5);
					\coordinate (O) at (1,-0.75); % Tâm đáy (giao AC và BD)
					\coordinate (S) at (1,3.5); % Đỉnh S thẳng đứng từ O
					
					% Vẽ các nét khuất
					\draw[dashed, thick] (A) -- (B) (A) -- (D) (A) -- (C) (B) -- (D) (S) -- (A) (S) -- (O);
					% Vẽ các nét thấy
					\draw[thick] (B) -- (C) -- (D) -- (S) -- (B) (S) -- (C);
					
					% Gắn nhãn
					\node[left] at (A) {$A$};
					\node[below left] at (B) {$B$};
					\node[below right] at (C) {$C$};
					\node[right] at (D) {$D$};
					\node[above] at (S) {$S$};
					\node[below, yshift=-1mm] at (O) {$O$};
					
					% Chấm điểm O
					\fill (O) circle (1.5pt);
				\end{tikzpicture}
			\end{center}
		\end{minipage}
		
	\end{enumerate}
	
	
	
	
	
	% ================================================
	% ĐỀ SỐ 2
	% ================================================
	\newpage
	\begin{center}
		{\LARGE \textbf{ĐỀ SỐ 2}}\\[4pt]
	\end{center}
	\vspace{4mm}
	
	% --- PHẦN 1: TRẮC NGHIỆM ---
	\noindent\textbf{PHẦN I. Câu trắc nghiệm nhiều phương án lựa chọn.} Thí sinh trả lời từ câu 1 đến câu 12. Mỗi câu hỏi thí sinh chỉ chọn một phương án.
	\vspace{2mm}
	
	\begin{enumerate}[label=\textbf{Câu \arabic*.}, leftmargin=*]
		
		% CÂU 1: Vector không tọa độ (Đảo hình hộp thành hình chóp tam giác/tứ diện)
		\item Cho tứ diện $ABCD$ có $G$ là trọng tâm của tứ diện. Mệnh đề nào dưới đây là đúng?
		\begin{multicols}{2}
			\begin{enumerate}[label=\Alph*.]
				\item $\overrightarrow{GA} + \overrightarrow{GB} + \overrightarrow{GC} + \overrightarrow{GD} = \vec{0}$.
				\item $\overrightarrow{GA} + \overrightarrow{GB} + \overrightarrow{GC} = \overrightarrow{GD}$.
				\item $\overrightarrow{GA} + \overrightarrow{GB} + \overrightarrow{GC} + \overrightarrow{GD} = 4\overrightarrow{AB}$.
				\item $\overrightarrow{AG} = \dfrac{1}{3}\left( \overrightarrow{AB} + \overrightarrow{AC} + \overrightarrow{AD} \right)$.
			\end{enumerate}
		\end{multicols}
		
		% CÂU 2: Hàm số - Đơn điệu (Đề 1 cho Đồ thị -> Đề 2 cho Công thức)
		\item Cho hàm số $y = -x^3 + 3x^2 + 9x - 1$. Hàm số đã cho đồng biến trên khoảng nào dưới đây?
		\begin{multicols}{4}
			\begin{enumerate}[label=\Alph*.]
				\item $(-\infty; -1)$.
				\item $(-1; 3)$.
				\item $(3; +\infty)$.
				\item $(-3; 1)$.
			\end{enumerate}
		\end{multicols}
		
		% CÂU 3: Thống kê (Đề 1 hỏi Tứ phân vị -> Đề 2 hỏi Khoảng biến thiên)
		\item Khảo sát độ tuổi của một nhóm người tham gia câu lạc bộ dưỡng sinh, ta thu được mẫu số liệu ghép nhóm sau:
		\begin{center}
			\begin{tabular}{|c|c|c|c|c|c|}
				\hline
				\textbf{Độ tuổi} & $[45; 50)$ & $[50; 55)$ & $[55; 60)$ & $[60; 65)$ & $[65; 70)$ \\
				\hline
				\textbf{Số người} & 5 & 12 & 18 & 10 & 5 \\
				\hline
			\end{tabular}
		\end{center}
		Tìm khoảng biến thiên $R$ của mẫu số liệu ghép nhóm trên.
		\begin{multicols}{4}
			\begin{enumerate}[label=\Alph*.]
				\item $R = 25$.
				\item $R = 30$.
				\item $R = 70$.
				\item $R = 45$.
			\end{enumerate}
		\end{multicols}
		
		% CÂU 4: Hàm số - Max/Min (Đề 1 cho Công thức -> Đề 2 cho Đồ thị)
		\item Cho hàm số $y = f(x)$ liên tục trên đoạn $[-2; 3]$ và có đồ thị như hình vẽ bên dưới. 
		\begin{center}
			\begin{tikzpicture}[scale=0.8]
				\begin{axis}[
					axis lines = middle,
					xlabel = {$x$}, ylabel = {$y$},
					ymin = -3, ymax = 5,
					xmin = -3, xmax = 4,
					xtick = {-2, -1, 1, 3}, ytick = {-2, -1, 4},
					ticklabel style = {fill=white, inner sep=1pt},
					width=8cm, height=6.5cm,
					samples=100,
					axis line style={->}
					]
					% Vẽ 1 đường cong đi qua các điểm (-2,0), (-1,-1), (1,4), (3,-2)
					\addplot [domain=-2:3, thick, blue, smooth] coordinates {(-2,0) (-1,-1) (0,2) (1,4) (2,2) (3,-2)};
					
					\draw[dashed, gray] (-1,0) -- (-1,-1) -- (0,-1);
					\draw[dashed, gray] (1,0) -- (1,4) -- (0,4);
					\draw[dashed, gray] (3,0) -- (3,-2) -- (0,-2);
					
					\fill (-2,0) circle (1.5pt);
					\fill (-1,-1) circle (1.5pt);
					\fill (1,4) circle (1.5pt);
					\fill (3,-2) circle (1.5pt);
					
					\node[below left] at (axis cs:0,0) {$O$};
				\end{axis}
			\end{tikzpicture}
		\end{center}
		Gọi $M, m$ lần lượt là giá trị lớn nhất và nhỏ nhất của hàm số trên đoạn $[-2; 3]$. Tính $M - m$.
		\begin{multicols}{4}
			\begin{enumerate}[label=\Alph*.]
				\item $6$.
				\item $5$.
				\item $4$.
				\item $2$.
			\end{enumerate}
		\end{multicols}
		
		% CÂU 5: Hàm số - Tiệm cận (Đề 1 cho Công thức -> Đề 2 cho BBT)
		\item Cho hàm số $y = f(x)$ có bảng biến thiên như sau:
		\begin{center}
			\begin{tikzpicture}[xscale=1, yscale=1]
				% Khung ngoài và các đường kẻ ngang dọc
				\draw (0,0) rectangle (8,3);
				\draw (0,1.5) -- (8,1.5); % Đường kẻ giữa y và y'
				\draw (0,2.2) -- (8,2.2); % Đường kẻ giữa y' và x
				\draw (1.2,0) -- (1.2,3); % Đường kẻ dọc ngăn cách
				
				% Dòng x
				\node at (0.6, 2.6) {$x$};
				\node at (1.8, 2.6) {$-\infty$};
				\node at (4.6, 2.6) {$1$};
				\node at (7.4, 2.6) {$+\infty$};
				
				% Dòng y'
				\node at (0.6, 1.85) {$y'$};
				\node at (3.2, 1.85) {$-$};
				\node at (6.0, 1.85) {$-$};
				
				% Nét đôi tại x = 1 (Không xác định)
				\draw (4.55, 0) -- (4.55, 2.2);
				\draw (4.65, 0) -- (4.65, 2.2);
				
				% Dòng y (Bố trí cao độ chuẩn xác)
				\node at (0.6, 0.75) {$y$};
				\node at (1.8, 1.15) {$2$};          % Trên cao
				\node at (4.2, 0.25) {$-\infty$};    % Dưới thấp (trái nét đôi)
				\node at (5.0, 1.15) {$+\infty$};    % Trên cao (phải nét đôi)
				\node at (7.4, 0.25) {$2$};          % Dưới thấp
				
				% Mũi tên xanh
				\draw[->, thick, blue] (2.1, 1.05) -- (3.9, 0.35);
				\draw[->, thick, blue] (5.3, 1.05) -- (7.1, 0.35);
			\end{tikzpicture}
		\end{center}
		Tổng số đường tiệm cận đứng và tiệm cận ngang của đồ thị hàm số đã cho là:
		\begin{multicols}{4}
			\begin{enumerate}[label=\Alph*.]
				\item $1$.
				\item $2$.
				\item $3$.
				\item $4$.
			\end{enumerate}
		\end{multicols}
		
		% CÂU 6: Vector Oxyz - Khái niệm cơ bản (Hình chiếu)
		\item Trong không gian $Oxyz$, cho điểm $M(3; -2; 5)$. Tọa độ hình chiếu vuông góc của $M$ lên mặt phẳng tọa độ $(Oyz)$ là:
		\begin{multicols}{4}
			\begin{enumerate}[label=\Alph*.]
				\item $M'(3; 0; 0)$.
				\item $M'(0; -2; 0)$.
				\item $M'(0; 0; 5)$.
				\item $M'(0; -2; 5)$.
			\end{enumerate}
		\end{multicols}
		
		% CÂU 7: Hàm số - Nhận diện đồ thị (Đề 1 cho Đồ thị -> Đề 2 cho BBT)
		\item Cho hàm số phân thức bậc nhất trên bậc nhất $y = \dfrac{ax+b}{cx+d}$ có bảng biến thiên như sau:
		\begin{center}
			\begin{tikzpicture}[xscale=1, yscale=1]
				% Khung ngoài và các đường kẻ ngang dọc
				\draw (0,0) rectangle (8,3);
				\draw (0,1.5) -- (8,1.5); % Đường kẻ giữa y và y'
				\draw (0,2.2) -- (8,2.2); % Đường kẻ giữa y' và x
				\draw (1.2,0) -- (1.2,3); % Đường kẻ dọc ngăn cách
				
				% Dòng x
				\node at (0.6, 2.6) {$x$};
				\node at (1.8, 2.6) {$-\infty$};
				\node at (4.6, 2.6) {$-1$};
				\node at (7.4, 2.6) {$+\infty$};
				
				% Dòng y'
				\node at (0.6, 1.85) {$y'$};
				\node at (3.2, 1.85) {$+$};
				\node at (6.0, 1.85) {$+$};
				
				% Nét đôi tại x = -1 (Không xác định)
				\draw (4.55, 0) -- (4.55, 2.2);
				\draw (4.65, 0) -- (4.65, 2.2);
				
				% Dòng y (Bố trí cao độ chuẩn xác)
				\node at (0.6, 0.75) {$y$};
				\node at (1.8, 0.25) {$2$};          % Dưới thấp
				\node at (4.2, 1.15) {$+\infty$};    % Trên cao (trái nét đôi)
				\node at (5.0, 0.25) {$-\infty$};    % Dưới thấp (phải nét đôi)
				\node at (7.4, 1.15) {$2$};          % Trên cao
				
				% Mũi tên xanh
				\draw[->, thick, blue] (2.1, 0.35) -- (3.9, 1.05);
				\draw[->, thick, blue] (5.3, 0.35) -- (7.1, 1.05);
			\end{tikzpicture}
		\end{center}
		Hàm số đã cho là hàm số nào dưới đây?
		\begin{multicols}{2}
			\begin{enumerate}[label=\Alph*.]
				\item $y = \dfrac{2x-1}{x+1}$.
				\item $y = \dfrac{2x+3}{x+1}$.
				\item $y = \dfrac{2x-1}{x-1}$.
				\item $y = \dfrac{-x+2}{x+1}$.
			\end{enumerate}
		\end{multicols}
		
		% CÂU 8: Hàm số - Max/Min (Đề 1 cho BBT -> Đề 2 cho Công thức)
		\item Tìm giá trị nhỏ nhất của hàm số $y = x^3 - 3x^2 - 9x + 5$ trên đoạn $[0; 4]$.
		\begin{multicols}{4}
			\begin{enumerate}[label=\Alph*.]
				\item $5$.
				\item $-22$.
				\item $-15$.
				\item $32$.
			\end{enumerate}
		\end{multicols}
		
		% CÂU 9: Vector Oxyz - Tính toán tọa độ (Hai vecto cùng phương)
		\item Trong không gian $Oxyz$, cho hai vecto $\vec{u} = (2; -1; 3)$ và $\vec{v} = (-4; m; -6)$. Tìm giá trị thực của tham số $m$ để hai vecto $\vec{u}$ và $\vec{v}$ cùng phương.
		\begin{multicols}{4}
			\begin{enumerate}[label=\Alph*.]
				\item $m = 2$.
				\item $m = -2$.
				\item $m = \dfrac{1}{2}$.
				\item $m = 4$.
			\end{enumerate}
		\end{multicols}
		
		% CÂU 10: Thống kê (Đề 1 cho Độ lệch chuẩn -> Đề 2 cho Phương sai)
		\item Khảo sát thời gian đọc sách (giờ/tuần) của 20 học sinh, ta có bảng số liệu ghép nhóm sau:
		\begin{center}
			\begin{tabular}{|c|c|c|c|}
				\hline
				\textbf{Thời gian (giờ)} & $[0; 2)$ & $[2; 4)$ & $[4; 6)$ \\
				\hline
				\textbf{Số học sinh} & 5 & 8 & 7 \\
				\hline
			\end{tabular}
		\end{center}
		Phương sai của mẫu số liệu ghép nhóm trên có giá trị bằng bao nhiêu (làm tròn đến hàng phần trăm)?
		\begin{multicols}{4}
			\begin{enumerate}[label=\Alph*.]
				\item $2,36$.
				\item $1,54$.
				\item $4,72$.
				\item $3,20$.
			\end{enumerate}
		\end{multicols}
		
		% CÂU 11: Vector Oxyz - Trọng tâm tam giác
		\item Trong không gian $Oxyz$, cho tam giác $ABC$ với tọa độ các đỉnh là $A(1; 2; -1)$, $B(2; -1; 3)$ và $C(-3; 5; 1)$. Tọa độ trọng tâm $G$ của tam giác $ABC$ là:
		\begin{multicols}{4}
			\begin{enumerate}[label=\Alph*.]
				\item $G(0; 6; 3)$.
				\item $G(0; 2; 1)$.
				\item $G(0; -2; -1)$.
				\item $G\left(0; 3; \dfrac{3}{2}\right)$.
			\end{enumerate}
		\end{multicols}
		
		% CÂU 12: Vector HHKG thuần túy - Tích vô hướng
		\item Cho hình vuông $ABCD$ có cạnh bằng $a$. Tính giá trị của biểu thức tích vô hướng $\overrightarrow{AB} \cdot \overrightarrow{AC}$.
		\begin{multicols}{4}
			\begin{enumerate}[label=\Alph*.]
				\item $a^2\sqrt{2}$.
				\item $2a^2$.
				\item $a^2$.
				\item $0$.
			\end{enumerate}
		\end{multicols}
		
	\end{enumerate}
	
	\vspace{4mm}
	% ======================================================
	% PHẦN 2: TRẮC NGHIỆM ĐÚNG SAI
	% ======================================================
	\noindent\textbf{PHẦN II. Câu trắc nghiệm đúng sai.} Thí sinh trả lời từ câu 1 đến câu 2. Trong mỗi ý a), b), c), d) ở mỗi câu, thí sinh chọn đúng hoặc sai.
	\vspace{2mm}
	
	\begin{enumerate}[label=\textbf{Câu \arabic*.}, leftmargin=*]
		
		% CÂU 1 ĐÚNG SAI: Khảo sát hàm số (3 ý cơ bản, 1 ý khó về tiệm cận xiên)
		\item Cho hàm số $y = \dfrac{x^2 - 3x + 6}{x - 1}$ có đồ thị là đường cong $(C)$. Xét tính đúng hay sai của các mệnh đề sau:
		\begin{enumerate}[label=\alph*)]
			\item Hàm số đã cho đạt cực tiểu tại điểm $x = 3$.
			\item Giao điểm của đồ thị $(C)$ với trục tung là điểm $M(0; -6)$.
			\item Đồ thị $(C)$ có một đường tiệm cận đứng là $x = 1$ và không có tiệm cận ngang.
			\item Đường tiệm cận xiên của đồ thị $(C)$ tạo với hai trục tọa độ một tam giác có diện tích bằng $4$.
		\end{enumerate}
		
		\vspace{2mm}
		% CÂU 2 ĐÚNG SAI: Vecto trong không gian (Hình chóp, tam giác vuông cân)
		\item Cho hình chóp $S.ABCD$ có đáy $ABCD$ là hình bình hành. Biết mặt bên $SAB$ là tam giác vuông cân tại $S$ và có cạnh đáy $AB = 2a$. Gọi $M$ là trung điểm của cạnh $AB$. 
		\begin{center}
			\begin{tikzpicture}[scale=0.9, line join=round, line cap=round]
				% Tọa độ các đỉnh
				\coordinate (B) at (0,0);
				\coordinate (C) at (4,-0.5);
				\coordinate (A) at (2,1.5);
				\coordinate (D) at (6,1);
				\coordinate (M) at (1,0.75); % Trung điểm AB
				\coordinate (S) at (1.5,3.5); % Đỉnh S (căn ke để SAB vuông tại S)
				
				% Vẽ nét khuất
				\draw[dashed, thick] (B) -- (A) (A) -- (D) (S) -- (A) (S) -- (M);
				\fill (M) circle (1.5pt);
				
				% Vẽ nét thấy
				\draw[thick] (S) -- (B) -- (C) -- (D) -- (S) -- (C);
				
				% Gắn nhãn
				\node[below left] at (B) {$B$};
				\node[below] at (C) {$C$};
				\node[above] at (S) {$S$};
				\node[right] at (D) {$D$};
				\node[right, yshift=2mm] at (A) {$A$};
				\node[below left] at (M) {$M$};
			\end{tikzpicture}
		\end{center}
		Xét tính đúng hay sai của các mệnh đề sau:
		\begin{enumerate}[label=\alph*)]
			\item $\overrightarrow{AB} \cdot \overrightarrow{SM} = 0$.
			\item $2\overrightarrow{AM} + \overrightarrow{AD} = \overrightarrow{AC}$.
			\item $\overrightarrow{AB} \cdot \overrightarrow{SA} = -2a^2$.
			\item $\overrightarrow{CD} \cdot \overrightarrow{SB} = -a^2$.
		\end{enumerate}
		
	\end{enumerate}
	
	\vspace{4mm}
	% ======================================================
	% PHẦN 3: TRẮC NGHIỆM TRẢ LỜI NGẮN
	% ======================================================
	\noindent\textbf{PHẦN III. Câu trắc nghiệm trả lời ngắn.} Thí sinh trả lời từ câu 1 đến câu 4.
	\vspace{2mm}
	
	\begin{enumerate}[label=\textbf{Câu \arabic*.}, leftmargin=*]
		
		% CÂU 1 TRẢ LỜI NGẮN: Max/Min hàm số trên đoạn
		\item Tính tổng giá trị lớn nhất và giá trị nhỏ nhất của hàm số $y = x - 2 + \dfrac{9}{x}$ trên đoạn $[1; 4]$.
		
		\vspace{2mm}
		% CÂU 2 TRẢ LỜI NGẮN: Thống kê (Tính tứ phân vị Q3 - 5 nhóm)
		\item Điểm kiểm tra môn Toán của 50 học sinh được thống kê bởi mẫu số liệu ghép nhóm sau:
		\begin{center}
			\renewcommand{\arraystretch}{1.2}
			\begin{tabular}{|c|c|c|c|c|c|}
				\hline
				\textbf{Điểm thi} & $[4; 5)$ & $[5; 6)$ & $[6; 7)$ & $[7; 8)$ & $[8; 9)$ \\
				\hline
				\textbf{Số học sinh} & 4 & 10 & 16 & 12 & 8 \\
				\hline
			\end{tabular}
		\end{center}
		Tính tứ phân vị thứ ba ($Q_3$) của mẫu số liệu ghép nhóm trên (viết kết quả dưới dạng số thập phân).
		
		\vspace{2mm}
		% CÂU 3 TRẢ LỜI NGẮN: Tính toán biểu thức vecto (Tích vô hướng)
		\item Trong không gian, cho hai vecto $\vec{u}$ và $\vec{v}$ thỏa mãn độ dài $|\vec{u}| = 3$, $|\vec{v}| = 4$ và $|\vec{u} + \vec{v}| = 5$. Tính giá trị của biểu thức: $P = (\vec{u} - \vec{v}) \cdot (2\vec{u} + 3\vec{v})$.
		
		\vspace{2mm}
		% CÂU 4 TRẢ LỜI NGẮN: Hệ trục Oxyz (Điểm chia đoạn thẳng)
		\item Trong không gian $Oxyz$, cho hai điểm $A(2; -1; 4)$ và $B(-4; 5; -2)$. Biết tọa độ của điểm $M(a; b; c)$ thỏa mãn hệ thức $2\overrightarrow{MA} + \overrightarrow{MB} = \vec{0}$. Tính giá trị của biểu thức $T = a + b - c$.
		
	\end{enumerate}
	
	% ======================================================
	% PHẦN 4: TỰ LUẬN
	% ======================================================
	\vspace{4mm}
	\noindent\textbf{PHẦN IV. Tự luận.} Thí sinh trình bày bài giải chi tiết từ bài 1 đến bài 2.
	\vspace{2mm}
	
	\begin{enumerate}[label=\textbf{Bài \arabic*.}, leftmargin=*]
		
		% BÀI 1: KHẢO SÁT HÀM SỐ (1.5 ĐIỂM)
		\item Cho hàm số $y = -x^3 + 3x - 1$ có đồ thị là đường cong $(C)$.
		\begin{enumerate}[label=\alph*)]
			\item  Khảo sát sự biến thiên (đồng biến, nghịch biến, cực trị) và xác định tọa độ tâm đối xứng của đồ thị hàm số đã cho.
			\item  Biết đường thẳng $d: y = 2x - 1$ cắt đồ thị $(C)$ tại ba điểm phân biệt $A, B, C$ (giả sử $x_A < x_B < x_C$). Tìm tọa độ các điểm $A, B, C$ và chứng minh rằng điểm $B$ là trung điểm của đoạn thẳng $AC$.
		\end{enumerate}
		
		\vspace{4mm}
		% BÀI 2: VECTO TRONG KHÔNG GIAN (1.5 ĐIỂM)
		\item 
		\begin{enumerate}[label=\alph*)]
			\item  Cho tứ diện $ABCD$. Gọi $I, J$ lần lượt là trung điểm của các cạnh $AB$ và $CD$. Chứng minh đẳng thức vecto sau:
			\[ \overrightarrow{AC} + \overrightarrow{BD} = 2\overrightarrow{IJ} \]
			\item Giả sử tứ diện $ABCD$ ở câu trên là tứ diện đều có tất cả các cạnh bằng $a$ (Hình minh họa bên cạnh). Tính theo $a$: 
			\[ \overrightarrow{AB} \cdot \overrightarrow{CD} \quad \text{và} \quad \overrightarrow{AB} \cdot \overrightarrow{AJ} \]
		\end{enumerate}
		\hfill
		
	\end{enumerate}
	
	% ================================================
	% ĐỀ SỐ 3
	% ================================================
	\newpage
	\begin{center}
		{\LARGE \textbf{ĐỀ SỐ 3}}\\[4pt]
	\end{center}
	\vspace{4mm}
	
	% --- PHẦN 1: TRẮC NGHIỆM ---
	\noindent\textbf{PHẦN I. Câu trắc nghiệm nhiều phương án lựa chọn.} Thí sinh trả lời từ câu 1 đến câu 12. Mỗi câu hỏi thí sinh chỉ chọn một phương án.
	\vspace{2mm}
	
	\begin{enumerate}[label=\textbf{Câu \arabic*.}, leftmargin=*]
		
		% CÂU 1: Tiệm cận - Công thức (có bẫy nghiệm tử mẫu)
		\item Giao điểm $I$ của đường tiệm cận đứng và đường tiệm cận ngang của đồ thị hàm số $y = \dfrac{x^2 - 3x + 2}{x^2 - 1}$ có tọa độ là:
		\begin{multicols}{4}
			\begin{enumerate}[label=\Alph*.]
				\item $I(1; 1)$.
				\item $I(-1; 1)$.
				\item $I(-1; -2)$.
				\item $I(1; -1)$.
			\end{enumerate}
		\end{multicols}
		
		% CÂU 2: Vector thuần túy - Lục giác đều
		\item Cho lục giác đều $ABCDEF$ có tâm $O$. Vecto tổng $\overrightarrow{AB} + \overrightarrow{AF}$ bằng vecto nào dưới đây?
		\begin{center}
			\begin{tikzpicture}[scale=0.8]
				% Lục giác đều
				\foreach \i/\label/\pos in {0/C/right, 60/B/above right, 120/A/above left, 180/F/left, 240/E/below left, 300/D/below right} {
					\coordinate (\label) at (\i:2);
					\node[\pos] at (\label) {$\label$};
				}
				\coordinate (O) at (0,0);
				\node[below, yshift=-1mm] at (O) {$O$};
				\draw[thick] (A) -- (B) -- (C) -- (D) -- (E) -- (F) -- cycle;
				\draw[dashed, gray] (A) -- (D) (B) -- (E) (C) -- (F);
				\fill (O) circle (1.5pt);
			\end{tikzpicture}
		\end{center}
		\begin{multicols}{4}
			\begin{enumerate}[label=\Alph*.]
				\item $\overrightarrow{AO}$.
				\item $\overrightarrow{AD}$.
				\item $\overrightarrow{OA}$.
				\item $\overrightarrow{FC}$.
			\end{enumerate}
		\end{multicols}
		
		% CÂU 3: Thống kê - Khoảng tứ phân vị
		\item Bảng dưới đây thống kê cự ly ném tạ (đơn vị: mét) của 40 vận động viên:
		\begin{center}
			\begin{tabular}{|c|c|c|c|c|}
				\hline
				\textbf{Cự ly (m)} & $[10; 12)$ & $[12; 14)$ & $[14; 16)$ & $[16; 18)$ \\
				\hline
				\textbf{Số VĐV} & 5 & 15 & 12 & 8 \\
				\hline
			\end{tabular}
		\end{center}
		Khoảng tứ phân vị $\Delta_Q$ của mẫu số liệu ghép nhóm này xấp xỉ bằng bao nhiêu (làm tròn đến hàng phần chục)?
		\begin{multicols}{4}
			\begin{enumerate}[label=\Alph*.]
				\item $2,5$.
				\item $3,2$.
				\item $3,0$.
				\item $2,8$.
			\end{enumerate}
		\end{multicols}
		
		% CÂU 4: Max/Min 1 - Bảng biến thiên
		\item Cho hàm số $y = f(x)$ liên tục trên đoạn $[-1; 4]$ và có bảng biến thiên như hình dưới:
		\begin{center}
			\begin{tikzpicture}[xscale=1.2, yscale=0.8]
				\draw (0,0) rectangle (8,3);
				\draw (0,1) -- (8,1);
				\draw (0,2) -- (8,2);
				\draw (1,0) -- (1,3);
				
				\node at (0.5, 2.5) {$x$};
				\node at (1.5, 2.5) {$-1$};
				\node at (3.5, 2.5) {$1$};
				\node at (5.5, 2.5) {$3$};
				\node at (7.5, 2.5) {$4$};
				
				\node at (0.5, 1.5) {$y'$};
				\node at (2.5, 1.5) {$+$};
				\node at (3.5, 1.5) {$0$};
				\node at (4.5, 1.5) {$-$};
				\node at (5.5, 1.5) {$0$};
				\node at (6.5, 1.5) {$+$};
				
				\node at (0.5, 0.5) {$y$};
				\node at (1.5, 0.2) {$-2$};
				\node at (3.5, 0.7) {$5$};
				\node at (5.5, 0.2) {$0$};
				\node at (7.5, 0.7) {$3$};
				
				\draw[->, thick, blue] (1.8, 0.3) -- (3.2, 0.6);
				\draw[->, thick, blue] (3.8, 0.6) -- (5.2, 0.3);
				\draw[->, thick, blue] (5.8, 0.3) -- (7.2, 0.6);
			\end{tikzpicture}
		\end{center}
		Gọi $M$ là giá trị lớn nhất của hàm số trên đoạn $[-1; 4]$. Khẳng định nào sau đây đúng?
		\begin{multicols}{4}
			\begin{enumerate}[label=\Alph*.]
				\item $M = 3$.
				\item $M = 5$.
				\item $M = 0$.
				\item $M = -2$.
			\end{enumerate}
		\end{multicols}
		
		% CÂU 5: Vector tọa độ - Góc giữa 2 vecto
		\item Trong không gian $Oxyz$, cho hai vecto $\vec{u} = (1; 1; 0)$ và $\vec{v} = (1; 0; 1)$. Tính góc giữa hai vecto $\vec{u}$ và $\vec{v}$.
		\begin{multicols}{4}
			\begin{enumerate}[label=\Alph*.]
				\item $90^\circ$.
				\item $45^\circ$.
				\item $60^\circ$.
				\item $30^\circ$.
			\end{enumerate}
		\end{multicols}
		
		% CÂU 6: Hàm số - Nhận diện đồ thị phân thức
		\item Đường cong trong hình vẽ bên dưới là đồ thị của hàm số nào?
		\begin{center}
			\begin{tikzpicture}[scale=0.8]
				\begin{axis}[
					axis lines = middle,
					xlabel = {$x$}, ylabel = {$y$},
					ymin = -2, ymax = 5,
					xmin = -4, xmax = 2,
					xtick = {-2, -1}, ytick = {1, 2},
					ticklabel style = {fill=white, inner sep=1pt},
					width=8cm, height=6.5cm,
					samples=100,
					axis line style={->}
					]
					% Tiệm cận
					\draw[dashed, red, thick] (axis cs:-1, -2) -- (axis cs:-1, 5);
					\draw[dashed, red, thick] (axis cs:-4, 2) -- (axis cs:2, 2);
					
					% Đồ thị y = (2x+1)/(x+1)
					\addplot [domain=-4:-1.2, thick, blue, smooth] {(2*x + 1)/(x + 1)};
					\addplot [domain=-0.8:2, thick, blue, smooth] {(2*x + 1)/(x + 1)};
					
					% Điểm cắt trục
					\fill (axis cs:0, 1) circle (1.5pt);
					\fill (axis cs:-0.5, 0) circle (1.5pt);
					
					\node[below left] at (axis cs:0,0) {$O$};
				\end{axis}
			\end{tikzpicture}
		\end{center}
		\begin{multicols}{2}
			\begin{enumerate}[label=\Alph*.]
				\item $y = \dfrac{x-1}{x+1}$.
				\item $y = \dfrac{2x+1}{x+1}$.
				\item $y = \dfrac{2x-1}{x-1}$.
				\item $y = \dfrac{x+1}{x-1}$.
			\end{enumerate}
		\end{multicols}
		
		% CÂU 7: Vector tọa độ - Điểm đối xứng
		\item Trong không gian $Oxyz$, cho điểm $A(1; 2; -1)$ và $B(3; 0; 2)$. Tìm tọa độ điểm $C$ sao cho $B$ là trung điểm của đoạn thẳng $AC$.
		\begin{multicols}{4}
			\begin{enumerate}[label=\Alph*.]
				\item $C(5; -2; 5)$.
				\item $C(2; 1; 0.5)$.
				\item $C(4; 2; 1)$.
				\item $C(-1; 4; -4)$.
			\end{enumerate}
		\end{multicols}
		
		% CÂU 8: Hàm số - Đơn điệu từ đồ thị đạo hàm f'(x) (Bị che nghiệm)
		\item Cho hàm số $y = f(x)$. Biết hàm số đạo hàm $y = f'(x)$ có đồ thị là một parabol cắt trục hoành tại hai điểm có hoành độ $x_1, x_2$ (với $x_1 < x_2$).
		\begin{center}
			\begin{tikzpicture}[scale=0.8]
				\begin{axis}[
					axis lines = middle,
					xlabel = {$x$}, ylabel = {$f'(x)$},
					ymin = -3, ymax = 4,
					xmin = -2.5, xmax = 4.5,
					xtick = \empty, % Ẩn số trên trục x
					ytick = \empty,
					width=8cm, height=6cm,
					samples=100,
					axis line style={->}
					]
					% Parabol f'(x) = 0.5(x+1)(x-3)
					\addplot [domain=-2.2:4.2, thick, blue, smooth] {0.5*(x+1)*(x-3)};
					
					\node[below left, xshift=-2mm] at (axis cs:0,0) {$O$};
					
					% Vẽ vạch tick nhỏ trên trục Ox
					\draw (axis cs:-1, 0.1) -- (axis cs:-1, -0.1);
					\draw (axis cs:3, 0.1) -- (axis cs:3, -0.1);
					
					
				\end{axis}
			\end{tikzpicture}
		\end{center}
		Hàm số $y = f(x)$ đồng biến trên khoảng nào dưới đây?
		\begin{multicols}{4}
			\begin{enumerate}[label=\Alph*.]
				\item $(x_1; x_2)$.
				\item $(-\infty; x_1)$.
				\item $(x_1; +\infty)$.
				\item $(0; x_2)$.
			\end{enumerate}
		\end{multicols}
		
		% CÂU 9: Vector tọa độ - Điểm cách đều 2 điểm (Điểm thuộc trục)
		\item Trong không gian $Oxyz$, cho hai điểm $A(2; 0; 1)$ và $B(0; 2; 3)$. Biết điểm $M$ nằm trên trục tung $Oy$ sao cho tam giác $MAB$ cân tại $M$. Tọa độ của điểm $M$ là:
		\begin{multicols}{4}
			\begin{enumerate}[label=\Alph*.]
				\item $M(0; -2; 0)$.
				\item $M(0; 1; 0)$.
				\item $M(0; 3; 0)$.
				\item $M(0; 2; 0)$.
			\end{enumerate}
		\end{multicols}
		
		% CÂU 10: Thống kê - Phương sai/Độ lệch chuẩn
		\item Đo chiều cao (cm) của một nhóm học sinh nam lớp 12, ta có bảng tần số ghép nhóm sau:
		\begin{center}
			\begin{tabular}{|c|c|c|c|}
				\hline
				\textbf{Chiều cao (cm)} & $[160; 165)$ & $[165; 170)$ & $[170; 175)$ \\
				\hline
				\textbf{Số học sinh} & 2 & 6 & 2 \\
				\hline
			\end{tabular}
		\end{center}
		Tính độ lệch chuẩn của mẫu số liệu ghép nhóm này (làm tròn đến 2 chữ số thập phân).
		\begin{multicols}{4}
			\begin{enumerate}[label=\Alph*.]
				\item $13,50$.
				\item $3,16$.
				\item $10,00$.
				\item $4,25$.
			\end{enumerate}
		\end{multicols}
		
		% CÂU 11: Vector hình học - Lăng trụ tam giác
		\item Cho hình lăng trụ tam giác $ABC.A'B'C'$. Đặt $\vec{a} = \overrightarrow{AB}$, $\vec{b} = \overrightarrow{AC}$ và $\vec{c} = \overrightarrow{AA'}$. Khẳng định nào sau đây là đúng?
		\begin{multicols}{2}
			\begin{enumerate}[label=\Alph*.]
				\item $\overrightarrow{BC'} = -\vec{a} + \vec{b} + \vec{c}$.
				\item $\overrightarrow{BC'} = \vec{a} - \vec{b} + \vec{c}$.
				\item $\overrightarrow{BC'} = \vec{a} + \vec{b} - \vec{c}$.
				\item $\overrightarrow{BC'} = -\vec{a} - \vec{b} + \vec{c}$.
			\end{enumerate}
		\end{multicols}
		
		% CÂU 12: Max/Min - Công thức
		\item Giá trị lớn nhất của hàm số $y = x^4 - 2x^2 + 3$ trên đoạn $[0; 2]$ bằng:
		\begin{multicols}{4}
			\begin{enumerate}[label=\Alph*.]
				\item $3$.
				\item $2$.
				\item $11$.
				\item $0$.
			\end{enumerate}
		\end{multicols}
		
	\end{enumerate}
	
	\vspace{4mm}
	% ======================================================
	% PHẦN 2: TRẮC NGHIỆM ĐÚNG SAI
	% ======================================================
	\noindent\textbf{PHẦN II. Câu trắc nghiệm đúng sai.} Thí sinh trả lời từ câu 1 đến câu 2. Trong mỗi ý a), b), c), d) ở mỗi câu, thí sinh chọn đúng hoặc sai.
	\vspace{2mm}
	
	\begin{enumerate}[label=\textbf{Câu \arabic*.}, leftmargin=*]
		
		% CÂU 1 ĐÚNG SAI: Khảo sát hàm số (Dựa vào Bảng biến thiên)
		\item Cho hàm số $y = f(x)$ liên tục trên $\mathbb{R}$ và có bảng biến thiên như sau:
		\begin{center}
			\begin{tikzpicture}[xscale=1.2, yscale=0.8]
				\draw (0,0) rectangle (8,3);
				\draw (0,1) -- (8,1);
				\draw (0,2) -- (8,2);
				\draw (1,0) -- (1,3);
				
				\node at (0.5, 2.5) {$x$};
				\node at (1.5, 2.5) {$-\infty$};
				\node at (3.5, 2.5) {$-1$};
				\node at (5.5, 2.5) {$2$};
				\node at (7.5, 2.5) {$+\infty$};
				
				\node at (0.5, 1.5) {$f'(x)$};
				\node at (2.5, 1.5) {$-$};
				\node at (3.5, 1.5) {$0$};
				\node at (4.5, 1.5) {$+$};
				\node at (5.5, 1.5) {$0$};
				\node at (6.5, 1.5) {$-$};
				
				\node at (0.5, 0.5) {$f(x)$};
				\node at (1.5, 0.7) {$+\infty$};
				\node at (3.5, 0.2) {$-3$};
				\node at (5.5, 0.7) {$2$};
				\node at (7.5, 0.2) {$-\infty$};
				
				\draw[->, thick, blue] (1.8, 0.6) -- (3.2, 0.3);
				\draw[->, thick, blue] (3.8, 0.3) -- (5.2, 0.6);
				\draw[->, thick, blue] (5.8, 0.6) -- (7.2, 0.3);
			\end{tikzpicture}
		\end{center}
		Xét tính đúng hay sai của các mệnh đề sau:
		\begin{enumerate}[label=\alph*)]
			\item Hàm số nghịch biến trên khoảng $(-1; 2)$.
			\item Giá trị cực tiểu của hàm số bằng $-3$.
			\item Phương trình $f(x) = 0$ có đúng 3 nghiệm thực phân biệt.
			\item Giá trị lớn nhất của hàm số trên đoạn $[-2; 3]$ bằng $2$.
		\end{enumerate}
		
		\vspace{2mm}
		% CÂU 2 ĐÚNG SAI: Vecto HHKG (Hình lập phương)
		\item Cho hình lập phương $ABCD.A'B'C'D'$ có cạnh bằng $a$. 
		\begin{center}
			\begin{tikzpicture}[scale=1, line join=round, line cap=round]
				% Định nghĩa tọa độ
				\coordinate (A) at (0,0);
				\coordinate (B) at (1.5,1.5);
				\coordinate (D) at (3,0);
				\coordinate (C) at (4.5,1.5);
				\coordinate (A') at (0,3);
				\coordinate (B') at (1.5,4.5);
				\coordinate (D') at (3,3);
				\coordinate (C') at (4.5,4.5);
				
				% Điểm M trên C'B'
				\coordinate (M) at (2.5,4.5); 
				
				% Nét khuất
				\draw[dashed, thick] (A) -- (B) (B) -- (C) (B) -- (B') (A) -- (M);
				
				% Nét thấy
				\draw[thick] (A) -- (D) -- (C) -- (C') -- (D') -- (A') -- cycle;
				\draw[thick] (A') -- (B') -- (C') (D) -- (D');
				
				% Điểm M
				\fill[blue] (M) circle (1.5pt);
				
				% Gắn nhãn
				\node[below left] at (A) {$A$};
				\node[below right] at (B) {$B$};
				\node[below right] at (C) {$C$};
				\node[below] at (D) {$D$};
				\node[left] at (A') {$A'$};
				\node[above left] at (B') {$B'$};
				\node[above right] at (C') {$C'$};
				\node[right] at (D') {$D'$};
				\node[above] at (M) {$M$};
			\end{tikzpicture}
		\end{center}
		Xét tính đúng hay sai của các mệnh đề sau:
		\begin{enumerate}[label=\alph*)]
			\item $\left| \overrightarrow{A'B'} - \overrightarrow{D'C} \right| = a$.
			\item $\overrightarrow{AD} = \overrightarrow{C'B'}$.
			\item $\overrightarrow{AA'} \cdot \overrightarrow{B'C'} = 0$.
			\item Gọi $M$ là điểm thỏa mãn $\overrightarrow{C'M} = 2\overrightarrow{MB'}$. Khi đó $\overrightarrow{AM} = \overrightarrow{AA'} + \overrightarrow{AD} + \dfrac{2}{3}\overrightarrow{AB}$.
		\end{enumerate}
		
	\end{enumerate}
	
	\vspace{4mm}
	% ======================================================
	% PHẦN 3: TRẮC NGHIỆM TRẢ LỜI NGẮN
	% ======================================================
	\noindent\textbf{PHẦN III. Câu trắc nghiệm trả lời ngắn.} Thí sinh trả lời từ câu 1 đến câu 4.
	\vspace{2mm}
	
	\begin{enumerate}[label=\textbf{Câu \arabic*.}, leftmargin=*]
		
		% CÂU 1 TRẢ LỜI NGẮN: Max/Min hàm số
		\item Cho hàm số $y = \dfrac{x^2 - x + 4}{x}$ có giá trị lớn nhất là $M$ và giá trị nhỏ nhất là $m$ trên đoạn $[1; 4]$. Tính giá trị của biểu thức $T = M + 2m$.
		
		\vspace{2mm}
		% CÂU 2 TRẢ LỜI NGẮN: Thống kê (Trung vị / Tứ phân vị)
		\item Thống kê khối lượng (đơn vị: gram) của 50 quả cam thu hoạch được trong một khu vườn, ta được bảng số liệu ghép nhóm sau:
		\begin{center}
			\renewcommand{\arraystretch}{1.2}
			\begin{tabular}{|c|c|c|c|c|c|}
				\hline
				\textbf{Khối lượng (g)} & $[150; 160)$ & $[160; 170)$ & $[170; 180)$ & $[180; 190)$ & $[190; 200)$ \\
				\hline
				\textbf{Số quả cam} & 4 & 10 & 20 & 12 & 4 \\
				\hline
			\end{tabular}
		\end{center}
		Tính số trung vị (chính là tứ phân vị thứ hai $Q_2$) của mẫu số liệu ghép nhóm trên (làm tròn kết quả đến hàng phần chục).
		
		\vspace{2mm}
		% CÂU 3 TRẢ LỜI NGẮN: Oxyz (Tìm tọa độ điểm)
		\item Trong không gian $Oxyz$, cho hai điểm $A(1; -2; 3)$ và $B(-5; 4; 0)$. Tìm tọa độ điểm $M(a; b; c)$ nằm trên đoạn thẳng $AB$ sao cho $AM = 2MB$. Tính giá trị của biểu thức $P = a^2 + b^2 + c^2$.
		
		\vspace{2mm}
		% CÂU 4 TRẢ LỜI NGẮN: Vector hình không gian (Góc giữa 2 vecto)
		\item Trong không gian, cho hai vecto $\vec{a}$ và $\vec{b}$ khác vecto-không, thỏa mãn $|\vec{a}| = 2$, $|\vec{b}| = 4$ và $|2\vec{a} - \vec{b}| = 4\sqrt{2}$. Tính số đo góc giữa hai vecto $\vec{a}$ và $\vec{b}$ (viết kết quả tính theo đơn vị độ).
		
	\end{enumerate}
	
	% ======================================================
	% PHẦN 4: TỰ LUẬN
	% ======================================================
	\vspace{4mm}
	\noindent\textbf{PHẦN IV. Tự luận.} Thí sinh trình bày bài giải chi tiết từ bài 1 đến bài 2.
	\vspace{2mm}
	
	\begin{enumerate}[label=\textbf{Bài \arabic*.}, leftmargin=*]
		
		% BÀI 1: KHẢO SÁT HÀM SỐ (1.5 ĐIỂM)
		\item  Cho hàm số $y = \dfrac{2x - 1}{x + 1}$ có đồ thị là đường cong $(C)$.
		\begin{enumerate}[label=\alph*)]
			\item Khảo sát sự biến thiên (đồng biến, nghịch biến, cực trị, tiệm cận) và xác định tọa độ tâm đối xứng của đồ thị hàm số đã cho.
			\item Gọi $A$ và $B$ là hai điểm nằm trên hai nhánh khác nhau của đồ thị $(C)$. Tìm giá trị nhỏ nhất của độ dài đoạn thẳng $AB$.
		\end{enumerate}
		
		\vspace{4mm}
		% BÀI 2: VECTO TRONG KHÔNG GIAN (1.5 ĐIỂM)
		\item 
		\begin{enumerate}[label=\alph*)]
			\item  Cho hình lăng trụ $ABC.A'B'C'$. Gọi $G$ là trọng tâm của tam giác $ABC$ và $M$ là một điểm bất kỳ trong không gian. Chứng minh đẳng thức vecto sau:
			\[ \overrightarrow{MA'} + \overrightarrow{MB'} + \overrightarrow{MC'} = 3\overrightarrow{MG} + 3\overrightarrow{AA'} \]
			\item  Giả sử hình lăng trụ $ABC.A'B'C'$ ở câu trên là lăng trụ tam giác đều có tất cả các cạnh đều bằng $a$ (Hình minh họa bên cạnh). Gọi $N$ là trung điểm của cạnh $BB'$. Tính tích vô hướng sau theo $a$: 
			\[ \overrightarrow{A'B} \cdot \overrightarrow{A'N} \]
		\end{enumerate}
		
		\hfill
		\begin{center}
			\begin{tikzpicture}[scale=0.9, line join=round, line cap=round]
				% Lăng trụ tam giác
				\coordinate (A) at (0,0);
				\coordinate (B) at (1.5,-1.5);
				\coordinate (C) at (4,0);
				\coordinate (A') at (0,3);
				\coordinate (B') at (1.5,1.5);
				\coordinate (C') at (4,3);
				
				% Vẽ nét khuất
				\draw[dashed, thick] (A) -- (C);
				
				% Vẽ nét thấy
				\draw[thick] (A) -- (B) -- (C) (A) -- (A') (B) -- (B') (C) -- (C') (A') -- (B') -- (C') -- (A');
				
				% Gắn nhãn
				\node[left] at (A) {$A$};
				\node[below] at (B) {$B$};
				\node[right] at (C) {$C$};
				\node[left] at (A') {$A'$};
				\node[above] at (B') {$B'$};
				\node[right] at (C') {$C'$};
			\end{tikzpicture}
		\end{center}
		
	\end{enumerate}
	
	% ================================================
	% ĐỀ SỐ 4
	% ================================================
	\newpage
	\begin{center}
		{\LARGE \textbf{ĐỀ SỐ 4}}\\[4pt]
	\end{center}
	\vspace{4mm}
	
	% --- PHẦN 1: TRẮC NGHIỆM ---
	\noindent\textbf{PHẦN I. Câu trắc nghiệm nhiều phương án lựa chọn.} Thí sinh trả lời từ câu 1 đến câu 12. Mỗi câu hỏi thí sinh chỉ chọn một phương án.
	\vspace{2mm}
	
	\begin{enumerate}[label=\textbf{Câu \arabic*.}, leftmargin=*]
		
		% CÂU 1: Hàm số - Tiệm cận từ đồ thị
		\item Cho hàm số $y = f(x)$ có đồ thị như hình vẽ bên dưới. Phương trình đường tiệm cận đứng của đồ thị hàm số là:
		\begin{center}
			\begin{tikzpicture}[scale=0.8]
				\begin{axis}[
					axis lines = middle,
					xlabel = {$x$}, ylabel = {$y$},
					ymin = -3, ymax = 5,
					xmin = -3, xmax = 5,
					xtick = {2}, ytick = {1},
					ticklabel style = {fill=white, inner sep=1pt},
					width=8cm, height=6.5cm,
					samples=100,
					axis line style={->}
					]
					% Tiệm cận
					\draw[dashed, red, thick] (axis cs:2, -3) -- (axis cs:2, 5);
					\draw[dashed, red, thick] (axis cs:-3, 1) -- (axis cs:5, 1);
					
					% Đồ thị y = (x+1)/(x-2)
					\addplot [domain=-3:1.7, thick, blue, smooth] {(x + 1)/(x - 2)};
					\addplot [domain=2.3:5, thick, blue, smooth] {(x + 1)/(x - 2)};
					
					\node[below left] at (axis cs:0,0) {$O$};
				\end{axis}
			\end{tikzpicture}
		\end{center}
		\begin{multicols}{4}
			\begin{enumerate}[label=\Alph*.]
				\item $x = 2$.
				\item $y = 1$.
				\item $x = 1$.
				\item $y = 2$.
			\end{enumerate}
		\end{multicols}
		
		% CÂU 2: Oxyz - Phép toán tọa độ
		\item Trong không gian $Oxyz$, cho hai vecto $\vec{a} = (1; -2; 3)$ và $\vec{b} = (0; 1; -1)$. Tọa độ của vecto $\vec{c} = 2\vec{a} - 3\vec{b}$ là:
		\begin{multicols}{4}
			\begin{enumerate}[label=\Alph*.]
				\item $\vec{c} = (2; -1; 2)$.
				\item $\vec{c} = (2; -7; 9)$.
				\item $\vec{c} = (1; -1; 2)$.
				\item $\vec{c} = (2; -5; 5)$.
			\end{enumerate}
		\end{multicols}
		
		% CÂU 3: Thống kê - Khoảng biến thiên
		\item Khảo sát điểm thi môn Tiếng Anh của một nhóm học sinh, ta có mẫu số liệu ghép nhóm sau:
		\begin{center}
			\begin{tabular}{|c|c|c|c|c|}
				\hline
				\textbf{Điểm thi} & $[4; 5)$ & $[5; 6)$ & $[6; 7)$ & $[7; 8)$ \\
				\hline
				\textbf{Số học sinh} & 5 & 12 & 18 & 5 \\
				\hline
			\end{tabular}
		\end{center}
		Khoảng biến thiên $R$ của mẫu số liệu ghép nhóm trên bằng:
		\begin{multicols}{4}
			\begin{enumerate}[label=\Alph*.]
				\item $R = 4$.
				\item $R = 3$.
				\item $R = 5$.
				\item $R = 8$.
			\end{enumerate}
		\end{multicols}
		
		% CÂU 4: Hàm số - Max/Min (Công thức)
		\item Gọi $M, m$ lần lượt là giá trị lớn nhất và giá trị nhỏ nhất của hàm số $y = \dfrac{2x - 1}{x + 1}$ trên đoạn $[0; 2]$. Tính tích $M \cdot m$.
		\begin{multicols}{4}
			\begin{enumerate}[label=\Alph*.]
				\item $1$.
				\item $-2$.
				\item $0$.
				\item $-1$.
			\end{enumerate}
		\end{multicols}
		
		% CÂU 5: Vector thuần túy - Hình bình hành
		\item Cho hình bình hành $ABCD$. Mệnh đề nào dưới đây là \textbf{đúng}?
		\begin{multicols}{2}
			\begin{enumerate}[label=\Alph*.]
				\item $\overrightarrow{AB} + \overrightarrow{BC} = \overrightarrow{CA}$.
				\item $\overrightarrow{AB} + \overrightarrow{AD} = \overrightarrow{AC}$.
				\item $\overrightarrow{AB} - \overrightarrow{AD} = \overrightarrow{AC}$.
				\item $\overrightarrow{OA} + \overrightarrow{OB} = \vec{0}$ (với $O$ là tâm hình bình hành).
			\end{enumerate}
		\end{multicols}
		
		% CÂU 6: Hàm số - Xác định hàm từ BBT
		\item Bảng biến thiên dưới đây là của hàm số nào?
		\begin{center}
			\begin{tikzpicture}[xscale=1.2, yscale=0.8]
				\draw (0,0) rectangle (8,3);
				\draw (0,1) -- (8,1);
				\draw (0,2) -- (8,2);
				\draw (1,0) -- (1,3);
				
				\node at (0.5, 2.5) {$x$};
				\node at (1.5, 2.5) {$-\infty$};
				\node at (3.5, 2.5) {$-1$};
				\node at (5.5, 2.5) {$1$};
				\node at (7.5, 2.5) {$+\infty$};
				
				\node at (0.5, 1.5) {$y'$};
				\node at (2.5, 1.5) {$+$};
				\node at (3.5, 1.5) {$0$};
				\node at (4.5, 1.5) {$-$};
				\node at (5.5, 1.5) {$0$};
				\node at (6.5, 1.5) {$+$};
				
				\node at (0.5, 0.5) {$y$};
				\node at (1.5, 0.2) {$-\infty$};
				\node at (3.5, 0.7) {$2$};
				\node at (5.5, 0.2) {$-2$};
				\node at (7.5, 0.7) {$+\infty$};
				
				\draw[->, thick, blue] (1.8, 0.3) -- (3.2, 0.6);
				\draw[->, thick, blue] (3.8, 0.6) -- (5.2, 0.3);
				\draw[->, thick, blue] (5.8, 0.3) -- (7.2, 0.6);
			\end{tikzpicture}
		\end{center}
		\begin{multicols}{2}
			\begin{enumerate}[label=\Alph*.]
				\item $y = -x^3 + 3x$.
				\item $y = x^3 - 3x^2$.
				\item $y = x^3 - 3x$.
				\item $y = \dfrac{2x-1}{x+1}$.
			\end{enumerate}
		\end{multicols}
		
		% CÂU 7: Oxyz - Điểm đối xứng
		\item Trong không gian $Oxyz$, cho điểm $A(1; -2; 4)$. Hình chiếu vuông góc của điểm $A$ trên mặt phẳng $(Oxy)$ có tọa độ là:
		\begin{multicols}{4}
			\begin{enumerate}[label=\Alph*.]
				\item $(1; -2; -4)$.
				\item $(0; 0; 4)$.
				\item $(1; -2; 0)$.
				\item $(1; 0; 0)$.
			\end{enumerate}
		\end{multicols}
		
		% CÂU 8: Hàm số - Cực trị từ đồ thị hàm số gốc
		\item Cho hàm số $y = f(x)$ liên tục trên $\mathbb{R}$ và có đồ thị như hình vẽ bên. 
		\begin{center}
			\begin{tikzpicture}[scale=0.8]
				\begin{axis}[
					axis lines = middle,
					xlabel = {$x$}, ylabel = {$y$},
					ymin = -3, ymax = 3,
					xmin = -2.5, xmax = 2.5,
					xtick = {-1, 1}, ytick = {-2, 2},
					ticklabel style = {fill=white, inner sep=1pt},
					width=8cm, height=6.5cm,
					samples=100,
					axis line style={->}
					]
					% Hàm trùng phương dạng W ngược: y = -x^4 + 2x^2
					\addplot [domain=-1.7:1.7, thick, blue, smooth] {-x^4 + 2*x^2 + 1};
					
					\node[below left] at (axis cs:0,0) {$O$};
				\end{axis}
			\end{tikzpicture}
		\end{center}
		Hỏi hàm số đã cho có bao nhiêu điểm cực trị?
		\begin{multicols}{4}
			\begin{enumerate}[label=\Alph*.]
				\item $2$.
				\item $3$.
				\item $1$.
				\item $4$.
			\end{enumerate}
		\end{multicols}
		
		% CÂU 9: Oxyz - Độ dài đoạn thẳng
		\item Trong không gian $Oxyz$, cho hai điểm $A(1; 0; -1)$ và $B(2; 2; 1)$. Độ dài đoạn thẳng $AB$ bằng:
		\begin{multicols}{4}
			\begin{enumerate}[label=\Alph*.]
				\item $3$.
				\item $9$.
				\item $\sqrt{5}$.
				\item $\sqrt{3}$.
			\end{enumerate}
		\end{multicols}
		
		% CÂU 10: Thống kê - Độ lệch chuẩn
		\item Đo thời gian (phút) giải một câu đố của 10 học sinh, ta thu được bảng số liệu ghép nhóm sau:
		\begin{center}
			\begin{tabular}{|c|c|c|c|}
				\hline
				\textbf{Thời gian (phút)} & $[0; 2)$ & $[2; 4)$ & $[4; 6)$ \\
				\hline
				\textbf{Số học sinh} & 3 & 4 & 3 \\
				\hline
			\end{tabular}
		\end{center}
		Độ lệch chuẩn của mẫu số liệu ghép nhóm này xấp xỉ bằng bao nhiêu (làm tròn đến hai chữ số thập phân)?
		\begin{multicols}{4}
			\begin{enumerate}[label=\Alph*.]
				\item $2,40$.
				\item $1,41$.
				\item $1,55$.
				\item $3,00$.
			\end{enumerate}
		\end{multicols}
		
		% CÂU 11: Vector thuần túy - Tam giác trọng tâm
		\item Cho tam giác $ABC$ có trọng tâm $G$. Khẳng định nào sau đây là \textbf{đúng}?
		\begin{multicols}{2}
			\begin{enumerate}[label=\Alph*.]
				\item $\overrightarrow{AG} = \dfrac{1}{3}\left( \overrightarrow{AB} + \overrightarrow{AC} \right)$.
				\item $\overrightarrow{AG} = \dfrac{1}{2}\left( \overrightarrow{AB} + \overrightarrow{AC} \right)$.
				\item $\overrightarrow{AG} = \dfrac{2}{3}\left( \overrightarrow{AB} + \overrightarrow{AC} \right)$.
				\item $\overrightarrow{GA} + \overrightarrow{GB} + \overrightarrow{GC} = 3\overrightarrow{AG}$.
			\end{enumerate}
		\end{multicols}
		
		% CÂU 12: Hàm số - Max/Min trên đoạn (từ đồ thị)
		\item Cho hàm số $y = f(x)$ có đồ thị trên đoạn $[-2; 2]$ như hình vẽ.
		\begin{center}
			\begin{tikzpicture}[scale=0.8]
				\begin{axis}[
					axis lines = middle,
					xlabel = {$x$}, ylabel = {$y$},
					ymin = -3, ymax = 5,
					xmin = -2.5, xmax = 2.5,
					xtick = {-2, -1, 1, 2}, ytick = {-2, 4},
					ticklabel style = {fill=white, inner sep=1pt},
					width=8cm, height=6.5cm,
					samples=100,
					axis line style={->}
					]
					% Hàm bậc 3
					\addplot [domain=-2:2, thick, blue, smooth] {-x^3 + 3*x + 2};
					
					\draw[dashed, gray] (1,0) -- (1,4) -- (0,4);
					\draw[dashed, gray] (-1,0) -- (-1,0) -- (0,0); % Cực tiểu tại y=0
					\draw[dashed, gray] (-2,0) -- (-2,4) -- (0,4);
					\draw[dashed, gray] (2,0) -- (2,-4); % Tại 2 là 0
					
					\fill (1,4) circle (1.5pt);
					\fill (-1,0) circle (1.5pt);
					\fill (-2,4) circle (1.5pt);
					\fill (2,0) circle (1.5pt);
					
					\node[below left] at (axis cs:0,0) {$O$};
				\end{axis}
			\end{tikzpicture}
		\end{center}
		Tổng giá trị lớn nhất và giá trị nhỏ nhất của hàm số trên đoạn $[-2; 2]$ bằng:
		\begin{multicols}{4}
			\begin{enumerate}[label=\Alph*.]
				\item $2$.
				\item $4$.
				\item $6$.
				\item $0$.
			\end{enumerate}
		\end{multicols}
		
	\end{enumerate}
	
	\vspace{4mm}
	% ======================================================
	% PHẦN 2: TRẮC NGHIỆM ĐÚNG SAI
	% ======================================================
	\noindent\textbf{PHẦN II. Câu trắc nghiệm đúng sai.} Thí sinh trả lời từ câu 1 đến câu 2. Trong mỗi ý a), b), c), d) ở mỗi câu, thí sinh chọn đúng hoặc sai.
	\vspace{2mm}
	
	\begin{enumerate}[label=\textbf{Câu \arabic*.}, leftmargin=*]
		
		% CÂU 1 ĐÚNG SAI: Khảo sát hàm số (Đồ thị hàm phân thức 1/1)
		\item Cho hàm số $y = f(x) = \dfrac{ax+b}{cx+d}$ có đồ thị là đường cong $(C)$ như hình vẽ bên dưới. 
		\begin{center}
			\begin{tikzpicture}[scale=0.9]
				\begin{axis}[
					axis lines = middle,
					xlabel = {$x$}, ylabel = {$y$},
					ymin = -2, ymax = 6,
					xmin = -4, xmax = 4,
					xtick = {-1, 1}, ytick = {2},
					ticklabel style = {fill=white, inner sep=1pt},
					width=9cm, height=7cm,
					samples=100,
					axis line style={->}
					]
					% Tiệm cận đứng x = 1, Tiệm cận ngang y = 2
					\draw[dashed, red, thick] (axis cs:1, -2) -- (axis cs:1, 6);
					\draw[dashed, red, thick] (axis cs:-4, 2) -- (axis cs:4, 2);
					
					% Đồ thị y = (2x+1)/(x-1)
					\addplot [domain=-4:0.8, thick, blue, smooth] {(2*x + 1)/(x - 1)};
					\addplot [domain=1.2:4, thick, blue, smooth] {(2*x + 1)/(x - 1)};
					
					% Điểm cắt
					\fill (axis cs:0, -1) circle (1.5pt);
					\node[below left] at (axis cs:0,-1) {$-1$};
					
					\node[below left] at (axis cs:0,0) {$O$};
				\end{axis}
			\end{tikzpicture}
		\end{center}
		Xét tính đúng hay sai của các mệnh đề sau:
		\begin{enumerate}[label=\alph*)]
			\item Hàm số đã cho nghịch biến trên $\mathbb{R} \setminus \{1\}$.
			\item Tâm đối xứng của đồ thị $(C)$ là điểm $I(1; 2)$.
			\item Giao điểm của đồ thị $(C)$ với trục hoành có hoành độ bằng $-0{,}5$.
			\item Phương trình $f(x) = m$ (với $m$ là tham số thực) có hai nghiệm phân biệt khi và chỉ khi $m \neq 2$.
		\end{enumerate}
		
		\vspace{4mm}
		% CÂU 2 ĐÚNG SAI: Tọa độ Oxyz gắn với hình không gian (Hình hộp)
		\item Trong không gian $Oxyz$, cho hình hộp chữ nhật $OABC.O'A'B'C'$ có ba đỉnh nằm trên ba trục tọa độ là $A(2; 0; 0)$, $C(0; 4; 0)$ và $O'(0; 0; 3)$. (Điểm $O$ trùng với gốc tọa độ).
		\begin{center}
			\begin{tikzpicture}[scale=0.8, line join=round, line cap=round]
				% Trục Oxyz
				\draw[->, thick] (0,0) -- (-2,-1.5) node[below] {$x$};
				\draw[->, thick] (0,0) -- (6,0) node[right] {$y$};
				\draw[->, thick] (0,0) -- (0,5) node[above] {$z$};
				
				% Tọa độ các đỉnh 
				\coordinate (O) at (0,0);
				\coordinate (A) at (-1.5,-1.125); % Trục x vẽ xiên
				\coordinate (C) at (4,0);
				\coordinate (B) at (2.5,-1.125);
				\coordinate (O') at (0,3);
				\coordinate (A') at (-1.5,1.875);
				\coordinate (C') at (4,3);
				\coordinate (B') at (2.5,1.875);
				
				% Vẽ nét khuất
				\draw[dashed, thick] (O) -- (A) (O) -- (C) (O) -- (O');
				
				% Vẽ nét thấy
				\draw[thick] (A) -- (B) -- (C) (A) -- (A') (B) -- (B') (C) -- (C') (O') -- (A') -- (B') -- (C') -- cycle;
				
				% Gắn nhãn
				\node[below left] at (O) {$O$};
				\node[left] at (A) {$A$};
				\node[below] at (B) {$B$};
				\node[below right] at (C) {$C$};
				\node[above left] at (O') {$O'$};
				\node[left] at (A') {$A'$};
				\node[below right] at (B') {$B'$};
				\node[above right] at (C') {$C'$};
			\end{tikzpicture}
		\end{center}
		Xét tính đúng hay sai của các mệnh đề sau:
		\begin{enumerate}[label=\alph*)]
			\item Tọa độ của điểm $B'$ là $(2; 4; 3)$.
			\item Vecto $\overrightarrow{OB'}$ và vecto $\overrightarrow{AC'}$ vuông góc với nhau.
			\item Trung điểm của đoạn thẳng $AB'$ nằm trên mặt phẳng tọa độ $(Oyz)$.
			\item Diện tích của tam giác $OA'C$ bằng $\dfrac{\sqrt{181}}{2}$.
		\end{enumerate}
		
	\end{enumerate}
	
	\vspace{4mm}
	% ======================================================
	% PHẦN 3: TRẮC NGHIỆM TRẢ LỜI NGẮN
	% ======================================================
	\noindent\textbf{PHẦN III. Câu trắc nghiệm trả lời ngắn.} Thí sinh trả lời từ câu 1 đến câu 4.
	\vspace{2mm}
	
	\begin{enumerate}[label=\textbf{Câu \arabic*.}, leftmargin=*]
		
		% CÂU 1 TRẢ LỜI NGẮN: Khảo sát hàm số (Tìm khoảng nghịch biến)
		\item Cho hàm số $y = -\dfrac{1}{3}x^3 + 2x^2 + 5x - 1$. Biết hàm số đồng biến trên khoảng $(a; b)$ với $a, b$ là các số thực và $b - a$ đạt giá trị lớn nhất. Tính giá trị của biểu thức $T = 2b - a$.
		
		\vspace{2mm}
		% CÂU 2 TRẢ LỜI NGẮN: Thống kê (Tính Tứ phân vị Q1)
		\item Thời gian (phút) hoàn thành một bài kiểm tra của 40 học sinh được ghi lại ở bảng số liệu ghép nhóm sau:
		\begin{center}
			\renewcommand{\arraystretch}{1.2}
			\begin{tabular}{|c|c|c|c|c|c|}
				\hline
				\textbf{Thời gian (phút)} & $[10; 15)$ & $[15; 20)$ & $[20; 25)$ & $[25; 30)$ & $[30; 35)$ \\
				\hline
				\textbf{Số học sinh} & 2 & 8 & 16 & 10 & 4 \\
				\hline
			\end{tabular}
		\end{center}
		Tính tứ phân vị thứ nhất ($Q_1$) của mẫu số liệu ghép nhóm trên (viết kết quả dưới dạng số thập phân).
		
		\vspace{2mm}
		% CÂU 3 TRẢ LỜI NGẮN: Vector trong không gian (Tính góc)
		\item Cho hình chóp $S.ABC$ có $SA = SB = SC = a$ và các góc $\widehat{ASB} = 60^\circ, \widehat{BSC} = 90^\circ, \widehat{CSA} = 120^\circ$. Tính số đo của góc tạo bởi hai vecto $\overrightarrow{SA}$ và $\overrightarrow{BC}$ (viết kết quả tính theo đơn vị độ).
		
		\vspace{2mm}
		% CÂU 4 TRẢ LỜI NGẮN: Tọa độ Oxyz (Hình chiếu, khoảng cách)
		\item Trong không gian $Oxyz$, cho hai điểm $A(1; -2; 3)$ và $B(4; 2; -1)$. Gọi $A', B'$ lần lượt là hình chiếu vuông góc của $A$ và $B$ lên mặt phẳng tọa độ $(Oxz)$. Tính bình phương độ dài đoạn thẳng $A'B'$.
		
	\end{enumerate}
	
	% ======================================================
	% PHẦN 4: TỰ LUẬN
	% ======================================================
	\vspace{4mm}
	\noindent\textbf{PHẦN IV. Tự luận.} Thí sinh trình bày bài giải chi tiết từ bài 1 đến bài 2.
	\vspace{2mm}
	
	\begin{enumerate}[label=\textbf{Bài \arabic*.}, leftmargin=*]
		
		% BÀI 1: KHẢO SÁT HÀM SỐ (1.5 ĐIỂM)
		\item Cho hàm số $y = x^3 - 3x^2 + 2$ có đồ thị là đường cong $(C)$.
		\begin{enumerate}[label=\alph*)]
			\item Khảo sát sự biến thiên (đồng biến, nghịch biến, cực trị) và xác định tọa độ tâm đối xứng của đồ thị hàm số đã cho.
			\item Viết phương trình đường thẳng đi qua 2 điểm cục trị của  và tính khoảng cách từ điểm $M(-1; 3)$ đến đường thẳng này.
		\end{enumerate}
		
		\vspace{4mm}
		% BÀI 2: VECTO TRONG KHÔNG GIAN (1.5 ĐIỂM)
		\item
		\begin{enumerate}[label=\alph*)]
			\item Cho hình hộp $ABCD.A'B'C'D'$. Gọi $I$ là trung điểm của đoạn thẳng $AB$. Chứng minh đẳng thức vecto sau:
			\[ \overrightarrow{A'C} + \overrightarrow{A'D} = 2\overrightarrow{A'I} + \overrightarrow{A'C'} \]
			\item Giả sử hình hộp ở câu trên là hình hộp chữ nhật có các kích thước $AB = a$, $AD = 2a$ và $AA' = 3a$ (Hình minh họa bên cạnh). Tính giá trị của biểu thức tích vô hướng sau theo $a$:
			\[ \overrightarrow{AC'} \cdot \overrightarrow{BD'} \]
		\end{enumerate}
		\hfill
		\begin{center}
			\begin{tikzpicture}[scale=0.9, line join=round, line cap=round]
				% Hình hộp chữ nhật
				\coordinate (A) at (0,0);
				\coordinate (B) at (3,-0.5);
				\coordinate (D) at (1.5,1.5);
				\coordinate (C) at (4.5,1);
				\coordinate (A') at (0,3);
				\coordinate (B') at (3,2.5);
				\coordinate (D') at (1.5,4.5);
				\coordinate (C') at (4.5,4);
				
				% Vẽ nét khuất
				\draw[dashed, thick] (A) -- (D) (D) -- (C) (D) -- (D');
				
				% Vẽ nét thấy
				\draw[thick] (A) -- (B) -- (C) (A) -- (A') (B) -- (B') (C) -- (C') (A') -- (B') -- (C') -- (D') -- (A');
				
				% Gắn nhãn
				\node[below left] at (A) {$A$};
				\node[below right] at (B) {$B$};
				\node[right] at (C) {$C$};
				\node[left] at (D) {$D$};
				\node[left] at (A') {$A'$};
				\node[below right] at (B') {$B'$};
				\node[right] at (C') {$C'$};
				\node[above] at (D') {$D'$};
			\end{tikzpicture}
		\end{center}
		
	\end{enumerate}

	% ================================================
	% ĐỀ SỐ 5
	% ================================================
	\newpage
	\begin{center}
		{\LARGE \textbf{ĐỀ SỐ 5}}\\[4pt]
	\end{center}
	\vspace{4mm}
	
	% --- PHẦN 1: TRẮC NGHIỆM ---
	\noindent\textbf{PHẦN I. Câu trắc nghiệm nhiều phương án lựa chọn.} Thí sinh trả lời từ câu 1 đến câu 12. Mỗi câu hỏi thí sinh chỉ chọn một phương án.
	\vspace{2mm}
	
	\begin{enumerate}[label=\textbf{Câu \arabic*.}, leftmargin=*]
		
		% CÂU 1: Hàm số - Nhận diện dấu đồ thị phân thức
		\item Cho hàm số $y = \dfrac{ax+b}{cx+d}$ có đồ thị như hình vẽ bên dưới. Khẳng định nào sau đây là \textbf{đúng}?
		\begin{center}
			\begin{tikzpicture}[scale=0.8]
				\begin{axis}[
					axis lines = middle,
					xlabel = {$x$}, ylabel = {$y$},
					ymin = -4, ymax = 4,
					xmin = -4, xmax = 4,
					xtick = \empty, ytick = \empty,
					width=8cm, height=6.5cm,
					samples=100,
					axis line style={->}
					]
					% Tiệm cận đứng x = 1 (dương), Tiệm cận ngang y = 1 (dương)
					\draw[dashed, red, thick] (axis cs:1, -4) -- (axis cs:1, 4);
					\draw[dashed, red, thick] (axis cs:-4, 1) -- (axis cs:4, 1);
					
					% Đồ thị y = (x+1)/(x-1) -> cắt Oy tại -1 (âm), Ox tại -1 (âm)
					\addplot [domain=-4:0.7, thick, blue, smooth] {(x + 1)/(x - 1)};
					\addplot [domain=1.3:4, thick, blue, smooth] {(x + 1)/(x - 1)};
					
					\node[below left] at (axis cs:0,0) {$O$};
				\end{axis}
			\end{tikzpicture}
		\end{center}
		\begin{multicols}{2}
			\begin{enumerate}[label=\Alph*.]
				\item $ad > 0$ và $bc < 0$.
				\item $ad < 0$ và $ab > 0$.
				\item $ac > 0$ và $bd > 0$.
				\item $cd < 0$ và $bd < 0$.
			\end{enumerate}
		\end{multicols}
		
		% CÂU 2: Oxyz - Các phép toán tọa độ
		\item Trong không gian $Oxyz$, cho hai điểm $A(1; 2; -1)$ và $B(0; 3; 1)$. Tọa độ của vecto $\overrightarrow{AB}$ là:
		\begin{multicols}{4}
			\begin{enumerate}[label=\Alph*.]
				\item $(-1; 1; 2)$.
				\item $(1; -1; -2)$.
				\item $(1; 5; 0)$.
				\item $(-1; 1; 0)$.
			\end{enumerate}
		\end{multicols}
		
		% CÂU 3: Thống kê - Tìm nhóm chứa trung vị
		\item Thống kê số giờ tự học trong một tuần của 40 học sinh, ta có bảng tần số ghép nhóm sau:
		\begin{center}
			\begin{tabular}{|c|c|c|c|c|}
				\hline
				\textbf{Thời gian (giờ)} & $[0; 5)$ & $[5; 10)$ & $[10; 15)$ & $[15; 20)$ \\
				\hline
				\textbf{Số học sinh} & 6 & 15 & 12 & 7 \\
				\hline
			\end{tabular}
		\end{center}
		Nhóm nào dưới đây chứa số trung vị (tứ phân vị thứ hai $Q_2$) của mẫu số liệu trên?
		\begin{multicols}{4}
			\begin{enumerate}[label=\Alph*.]
				\item $[0; 5)$.
				\item $[5; 10)$.
				\item $[10; 15)$.
				\item $[15; 20)$.
			\end{enumerate}
		\end{multicols}
		
		% CÂU 4: Hàm số - Đơn điệu (Công thức)
		\item Cho hàm số $y = x^3 - 3x^2 + 2$. Hàm số đã cho nghịch biến trên khoảng nào dưới đây?
		\begin{multicols}{4}
			\begin{enumerate}[label=\Alph*.]
				\item $(-\infty; 0)$.
				\item $(2; +\infty)$.
				\item $(0; 2)$.
				\item $(-\infty; 2)$.
			\end{enumerate}
		\end{multicols}
		
		% CÂU 5: Vector thuần túy - Quy tắc hình hộp
		\item Cho hình hộp $ABCD.A'B'C'D'$. Đẳng thức vecto nào sau đây là \textbf{đúng}?
		\begin{multicols}{2}
			\begin{enumerate}[label=\Alph*.]
				\item $\overrightarrow{AC'} = \overrightarrow{AB} + \overrightarrow{AD} + \overrightarrow{AA'}$.
				\item $\overrightarrow{AC'} = \overrightarrow{AB} + \overrightarrow{BC} + \overrightarrow{CD}$.
				\item $\overrightarrow{AC'} = \overrightarrow{AD} + \overrightarrow{D'C'} + \overrightarrow{A'A}$.
				\item $\overrightarrow{AC'} = \overrightarrow{AB} + \overrightarrow{AD} - \overrightarrow{AA'}$.
			\end{enumerate}
		\end{multicols}
		
		% CÂU 6: Hàm số - Tiệm cận xiên
		\item Đường tiệm cận xiên của đồ thị hàm số $y = \dfrac{x^2 + 2x - 1}{x - 1}$ có phương trình là:
		\begin{multicols}{4}
			\begin{enumerate}[label=\Alph*.]
				\item $y = x - 1$.
				\item $y = x + 3$.
				\item $y = x + 2$.
				\item $y = x - 3$.
			\end{enumerate}
		\end{multicols}
		
		% CÂU 7: Oxyz - Khoảng cách đến mặt phẳng tọa độ
		\item Trong không gian $Oxyz$, khoảng cách từ điểm $M(3; -4; 5)$ đến mặt phẳng tọa độ $(Oxy)$ bằng:
		\begin{multicols}{4}
			\begin{enumerate}[label=\Alph*.]
				\item $3$.
				\item $4$.
				\item $5$.
				\item $\sqrt{34}$.
			\end{enumerate}
		\end{multicols}
		
		% CÂU 8: Hàm số - Max/Min từ Bảng biến thiên
		\item Cho hàm số $y = f(x)$ liên tục trên đoạn $[-2; 3]$ và có bảng biến thiên như sau:
		\begin{center}
			\begin{tikzpicture}[xscale=1.2, yscale=0.8]
				\draw (0,0) rectangle (8,3);
				\draw (0,1) -- (8,1);
				\draw (0,2) -- (8,2);
				\draw (1,0) -- (1,3);
				
				\node at (0.5, 2.5) {$x$};
				\node at (1.5, 2.5) {$-2$};
				\node at (3.5, 2.5) {$0$};
				\node at (5.5, 2.5) {$2$};
				\node at (7.5, 2.5) {$3$};
				
				\node at (0.5, 1.5) {$y'$};
				\node at (2.5, 1.5) {$+$};
				\node at (3.5, 1.5) {$0$};
				\node at (4.5, 1.5) {$-$};
				\node at (5.5, 1.5) {$0$};
				\node at (6.5, 1.5) {$+$};
				
				\node at (0.5, 0.5) {$y$};
				\node at (1.5, 0.2) {$1$};
				\node at (3.5, 0.7) {$4$};
				\node at (5.5, 0.2) {$-1$};
				\node at (7.5, 0.7) {$2$};
				
				\draw[->, thick, blue] (1.8, 0.3) -- (3.2, 0.6);
				\draw[->, thick, blue] (3.8, 0.6) -- (5.2, 0.3);
				\draw[->, thick, blue] (5.8, 0.3) -- (7.2, 0.6);
			\end{tikzpicture}
		\end{center}
		Gọi $M, m$ lần lượt là giá trị lớn nhất và giá trị nhỏ nhất của hàm số trên đoạn $[-2; 3]$. Tính $M + m$.
		\begin{multicols}{4}
			\begin{enumerate}[label=\Alph*.]
				\item $3$.
				\item $5$.
				\item $0$.
				\item $6$.
			\end{enumerate}
		\end{multicols}
		
		% CÂU 9: Vector thuần túy - Tích vô hướng (Bẫy góc tù)
		\item Cho tam giác đều $ABC$ có độ dài cạnh bằng $a$. Tính tích vô hướng $\overrightarrow{AB} \cdot \overrightarrow{BC}$.
		\begin{multicols}{4}
			\begin{enumerate}[label=\Alph*.]
				\item $\dfrac{a^2}{2}$.
				\item $-\dfrac{a^2}{2}$.
				\item $\dfrac{a^2\sqrt{3}}{2}$.
				\item $-\dfrac{a^2\sqrt{3}}{2}$.
			\end{enumerate}
		\end{multicols}
		
		% CÂU 10: Thống kê - Phương sai
		\item Doanh thu (triệu đồng) trong 4 ngày bán hàng của một cửa hàng được thống kê ở bảng sau:
		\begin{center}
			\begin{tabular}{|c|c|c|c|}
				\hline
				\textbf{Doanh thu} & $[0; 2)$ & $[2; 4)$ & $[4; 6)$ \\
				\hline
				\textbf{Số ngày} & 1 & 2 & 1 \\
				\hline
			\end{tabular}
		\end{center}
		Phương sai của mẫu số liệu ghép nhóm trên bằng bao nhiêu?
		\begin{multicols}{4}
			\begin{enumerate}[label=\Alph*.]
				\item $3$.
				\item $4$.
				\item $\sqrt{2}$.
				\item $2$.
			\end{enumerate}
		\end{multicols}
		
		% CÂU 11: Oxyz - Điểm đối xứng qua tâm
		\item Trong không gian $Oxyz$, cho điểm $M(2; -1; 3)$ và điểm $I(1; 1; 1)$. Tọa độ điểm $M'$ đối xứng với $M$ qua tâm $I$ là:
		\begin{multicols}{4}
			\begin{enumerate}[label=\Alph*.]
				\item $M'(0; 3; -1)$.
				\item $M'(1; 2; -2)$.
				\item $M'\left( \dfrac{3}{2}; 0; 2 \right)$.
				\item $M'(3; -3; 5)$.
			\end{enumerate}
		\end{multicols}
		
		% CÂU 12: Hàm số - Max/Min (Công thức)
		\item Tìm giá trị nhỏ nhất của hàm số $y = \dfrac{x - 1}{x + 2}$ trên đoạn $[0; 2]$.
		\begin{multicols}{4}
			\begin{enumerate}[label=\Alph*.]
				\item $-\dfrac{1}{2}$.
				\item $\dfrac{1}{4}$.
				\item $-1$.
				\item $0$.
			\end{enumerate}
		\end{multicols}
		
	\end{enumerate}
	
	\vspace{4mm}
	% ======================================================
	% PHẦN 2: TRẮC NGHIỆM ĐÚNG SAI
	% ======================================================
	\noindent\textbf{PHẦN II. Câu trắc nghiệm đúng sai.} Thí sinh trả lời từ câu 1 đến câu 2. Trong mỗi ý a), b), c), d) ở mỗi câu, thí sinh chọn đúng hoặc sai.
	\vspace{2mm}
	
	\begin{enumerate}[label=\textbf{Câu \arabic*.}, leftmargin=*]
		
		% CÂU 1 ĐÚNG SAI: Khảo sát hàm số (Đồ thị hàm bậc 3, tiếp tuyến, cực trị hàm hợp)
		\item Cho hàm số $y = f(x) = x^3 - 3x^2 + 1$ có đồ thị là đường cong $(C)$ như hình vẽ bên dưới.
		\begin{center}
			\begin{tikzpicture}[scale=0.9]
				\begin{axis}[
					axis lines = middle,
					xlabel = {$x$}, ylabel = {$y$},
					ymin = -4, ymax = 2.5,
					xmin = -1.5, xmax = 3.5,
					xtick = {1, 2}, ytick = {-3, 1},
					ticklabel style = {fill=white, inner sep=1pt},
					width=8cm, height=6.5cm,
					samples=100,
					axis line style={->}
					]
					% Đồ thị y = x^3 - 3x^2 + 1
					\addplot [domain=-1.2:3.2, thick, blue, smooth] {x^3 - 3*x^2 + 1};
					
					% Gióng tọa độ cực trị
					\draw[dashed, gray] (2,0) -- (2,-3) -- (0,-3);
					\fill (2,-3) circle (1.5pt);
					\fill (0,1) circle (1.5pt);
					
					% Điểm uốn
					\fill (1,-1) circle (1.5pt);
					\draw[dashed, gray] (1,0) -- (1,-1) -- (0,-1);
					\node[left] at (axis cs:0,-1) {$-1$};
					
					\node[below left] at (axis cs:0,0) {$O$};
				\end{axis}
			\end{tikzpicture}
		\end{center}
		Xét tính đúng hay sai của các mệnh đề sau:
		\begin{enumerate}[label=\alph*)]
			\item Tại điểm có hoành độ $x = 1$, đạo hàm của hàm số đạt giá trị bằng $0$.
			\item Đồ thị $(C)$ có điểm uốn là $I(1; -1)$.
			\item Phương trình tiếp tuyến của đồ thị $(C)$ tại điểm uốn là $y = -3x + 2$.
			\item Hàm số $g(x) = f(x^2 - 2x)$ có đúng $3$ điểm cực trị.
		\end{enumerate}
		
		\vspace{4mm}
		% CÂU 2 ĐÚNG SAI: Tọa độ Oxyz (Tam giác, trực tâm, diện tích, thể tích)
		\item Trong không gian $Oxyz$, cho ba điểm $A(1; 0; 0)$, $B(0; 2; 0)$ và $C(0; 0; 3)$. Xét tính đúng hay sai của các mệnh đề sau:
		\begin{enumerate}[label=\alph*)]
			\item Tích vô hướng $\overrightarrow{AB} \cdot \overrightarrow{AC} = 1$.
			\item Diện tích của tam giác $ABC$ bằng $\dfrac{7}{2}$.
			\item Gọi $H$ là trực tâm của tam giác $ABC$. Tọa độ của điểm $H$ là $H\left( \dfrac{36}{49}; \dfrac{18}{49}; \dfrac{12}{49} \right)$.
			\item Gọi $M(x; y; 0)$ là một điểm thuộc mặt phẳng $(Oxy)$. Tồn tại vô số điểm $M$ sao cho thể tích khối tứ diện $MABC$ bằng $2$.
		\end{enumerate}
		
	\end{enumerate}
	
	\vspace{4mm}
	% ======================================================
	% PHẦN 3: TRẮC NGHIỆM TRẢ LỜI NGẮN
	% ======================================================
	\noindent\textbf{PHẦN III. Câu trắc nghiệm trả lời ngắn.} Thí sinh trả lời từ câu 1 đến câu 4.
	\vspace{2mm}
	
	\begin{enumerate}[label=\textbf{Câu \arabic*.}, leftmargin=*]
		
		% CÂU 1 TRẢ LỜI NGẮN: Điểm cực trị hàm phân thức
		\item Biết đồ thị hàm số $y = \dfrac{x^2 + 2x + 5}{x + 1}$ có điểm cực tiểu là $M(x_0; y_0)$. Tính giá trị của tổng $S = x_0 + y_0$.
		
		\vspace{2mm}
		% CÂU 2 TRẢ LỜI NGẮN: Thống kê (Phương sai)
		\item Thời gian chờ (phút) của khách hàng tại một quầy thanh toán được ghi nhận trong bảng số liệu ghép nhóm sau:
		\begin{center}
			\renewcommand{\arraystretch}{1.2}
			\begin{tabular}{|c|c|c|c|}
				\hline
				\textbf{Thời gian (phút)} & $[0; 10)$ & $[10; 20)$ & $[20; 30)$ \\
				\hline
				\textbf{Số khách hàng} & 2 & 6 & 2 \\
				\hline
			\end{tabular}
		\end{center}
		Tính phương sai của mẫu số liệu ghép nhóm trên.
		
		\vspace{2mm}
		% CÂU 3 TRẢ LỜI NGẮN: Vector thuần túy (Độ dài tổng vecto)
		\item Cho hình vuông $ABCD$ có độ dài cạnh bằng $3$. Tính bình phương độ dài của vecto tổng: $\vec{u} = \overrightarrow{AB} + \overrightarrow{AC} + \overrightarrow{AD}$.
		
		\vspace{2mm}
		% CÂU 4 TRẢ LỜI NGẮN: Oxyz (Điểm cách đều)
		\item Trong không gian $Oxyz$, cho hai điểm $A(1; 2; 3)$ và $B(-1; 4; 1)$. Biết điểm $M$ nằm trên trục cao độ $Oz$ sao cho $MA = MB$. Tính cao độ của điểm $M$.
		
	\end{enumerate}
	
	% ======================================================
	% PHẦN 4: TỰ LUẬN
	% ======================================================
	\vspace{4mm}
	\noindent\textbf{PHẦN IV. Tự luận.} Thí sinh trình bày bài giải chi tiết từ bài 1 đến bài 2.
	\vspace{2mm}
	
	\begin{enumerate}[label=\textbf{Bài \arabic*.}, leftmargin=*]
		
		% BÀI 1: KHẢO SÁT HÀM SỐ (1.5 ĐIỂM)
		\item Cho hàm số $y = \dfrac{x^2 - x + 1}{x - 1}$ có đồ thị là đường cong $(C)$.
		\begin{enumerate}[label=\alph*)]
			\item Khảo sát sự biến thiên (đồng biến, nghịch biến, cực trị, tiệm cận) và xác định tọa độ tâm đối xứng của đồ thị hàm số đã cho.
			\item Tìm tất cả các điểm nằm trên đồ thị $(C)$ có tọa độ (hoành độ và tung độ) đều là các số nguyên.
		\end{enumerate}
		
		\vspace{4mm}
		% BÀI 2: VECTO TRONG KHÔNG GIAN (1.5 ĐIỂM)
		\item
		\begin{enumerate}[label=\alph*)]
			\item Cho hình lập phương $ABCD.A'B'C'D'$. Sử dụng các quy tắc cộng vecto, chứng minh đẳng thức:
			\[ \overrightarrow{AC'} = \overrightarrow{AB} + \overrightarrow{AD} + \overrightarrow{AA'} \]
			\item Giả sử hình lập phương ở câu trên có cạnh bằng $a$ (Hình minh họa dưới). Phân tích các vecto $\overrightarrow{AC'}$ và $\overrightarrow{A'B}$ theo 3 vecto cơ sở $\overrightarrow{AB}, \overrightarrow{AD}, \overrightarrow{AA'}$, từ đó tính tích vô hướng $\overrightarrow{AC'} \cdot \overrightarrow{A'B}$.
		\end{enumerate}
		\hfill
		\begin{center}
			\begin{tikzpicture}[scale=1, line join=round, line cap=round]
				% Hình lập phương
				\coordinate (A) at (0,0);
				\coordinate (B) at (1.5,-0.5);
				\coordinate (D) at (2,1);
				\coordinate (C) at (3.5,0.5);
				\coordinate (A') at (0,2.5);
				\coordinate (B') at (1.5,2);
				\coordinate (D') at (2,3.5);
				\coordinate (C') at (3.5,3);
				
				% Vẽ nét khuất
				\draw[dashed, thick] (A) -- (D) (D) -- (C) (D) -- (D');
				
				% Vẽ nét thấy
				\draw[thick] (A) -- (B) -- (C) (A) -- (A') (B) -- (B') (C) -- (C') (A') -- (B') -- (C') -- (D') -- (A');
				
				% Gắn nhãn
				\node[below left] at (A) {$A$};
				\node[below] at (B) {$B$};
				\node[right] at (C) {$C$};
				\node[above left] at (D) {$D$};
				\node[left] at (A') {$A'$};
				\node[below right] at (B') {$B'$};
				\node[right] at (C') {$C'$};
				\node[above] at (D') {$D'$};
			\end{tikzpicture}
		\end{center}
		
	\end{enumerate}

	
\end{document}